%% (H1) Společná informatika
\subsection{1. Automaty a jazyky}
\begin{itemize}
\item Regulární jazyky
\begin{itemize}
\item deterministický konečný automat (DFA) je pětice $A=(Q,\Sigma,\delta,q_0,F)$
\begin{itemize}
\item $Q$ \dots{} konečná množina stavů
\item $\Sigma$ \dots{} konečná neprázdná množina vstupních symbolů (abeceda)
\item $\delta:Q\times\Sigma\to Q$ \dots{} přechodová funkce
\item $q_0\in Q$ \dots{} počáteční stav
\item $F\subseteq Q$ \dots{} množina koncových (přijímajících) stavů
\end{itemize}
\item jazykem rozpoznávaným (přijímaným) deterministickým konečným automatem $A=(Q,\Sigma,\delta,q_0,F)$ nazveme jazyk $L(A)=\set{w\mid w\in\Sigma^*\land\delta^*(q_0,w)\in F}$
\item slovo je přijímáno automatem $A\equiv w\in L(A)$
\item jazyk $L$ je rozpoznatelný konečným automatem $\equiv$ existuje konečný automat $A$ takový, že $L=L(A)$
\item třídu jazyků rozpoznatelných konečnými automaty označíme $\mathcal F$, nazveme regulární jazyky
\item iterační (pumping) lemma
\begin{itemize}
\item nechť $L$ je regulární jazyk
\item $(\exists n\in\mathbb N)(\forall w\in L):|w|\geq n\implies w=xyz$
\begin{itemize}
\item $y\neq\epsilon$
\item $|xy|\leq n$
\item $\forall k\geq 0:xy^kz\in L$
\end{itemize}
\end{itemize}
\item věta: neregularita $L_{01}=\set{0^i1^i\mid i\geq 0}$
\begin{itemize}
\item předpokládejme regularitu $L_{01}$
\item vezměme $n$ z pumping lemmatu
\item zvolme $w=0^n1^n$
\item rozdělme $w=xyz$ dle pumping lemmatu
\begin{itemize}
\item $|xy|\leq n$ je na začátku $w$, takže obsahuje jen nuly
\item $y\neq\epsilon$
\end{itemize}
\item z pumping lemmatu $xz\in L_{01}$, což je spor, protože $xz$ obsahuje méně nul než jedniček
\end{itemize}
\item kongruence
\begin{itemize}
\item mějme konečnou abecedu $\Sigma$ a relaci ekvivalence $\sim$ na $\Sigma^*$
\item $\sim$ je pravá kongruence $\equiv(\forall u,v,w\in\Sigma^*):u\sim v\implies uw\sim vw$
\item $\sim$ je konečného indexu $\equiv$ rozklad $\Sigma^*/\sim$ má konečný počet tříd
\item třídu kongruence $\sim$ obsahující slovo $u$ značíme $[u]_\sim$ nebo $[u]$
\item příklady
\begin{itemize}
\item relace \quotedblbase{}končí stejným písmenem\textquotedblleft{} je pravá kongruence
\item relace \quotedblbase{}mají stejný počet znaků\textquotedblleft{} není konečného indexu
\end{itemize}
\end{itemize}
\item Myhill-Nerodova věta
\begin{itemize}
\item nechť $L$ je jazyk nad konečnou abecedou $\Sigma$
\item potom následující trzení jsou ekvivalentní:
\begin{itemize}
\item $L$ je rozpoznatelný konečným automatem
\item existuje pravá kongruence $\sim$ konečného indexu nad $\Sigma^*$ tak, že $L$ je sjednocením jistých tříd rozkladu $\Sigma^*/\sim$
\end{itemize}
\item použití k důkazu neregularity
\begin{itemize}
\item ukážeme pro pumpovatelný jazyk slov ve tvaru $a^+b^ic^i$ nebo $b^ic^j$
\item pro regulární $L$ musí existovat pravá kongruence konečného indexu $m$
\item mějme množinu řetězců $S=\set{ab,abb,abbb,\dots,ab^{m+1}}$
\item $|S|=m+1$, tedy nutně existují dvě slova $i\neq j$, která padnou do stejné třídy
\begin{itemize}
\item $ab^i\sim ab^j$
\item tedy $ab^ic^i\sim ab^jc^i$
\item ale $ab^ic^i\in L$ a $ab^jc^i\notin L$, což je spor, takže jazyk není regulární
\end{itemize}
\end{itemize}
\end{itemize}
\item věta: jazyk $L$ je rozpoznatelný $\epsilon$NFA, právě když je $L$ regulární
\begin{itemize}
\item pro libovolný $\epsilon$NFA $N$ zkonstruujeme DFA $D$ přijímající stejný jazyk jako $N$
\item nové stavy jsou $\epsilon$-uzavřené podmnožiny $Q_N$
\end{itemize}
\end{itemize}
\item Deterministický a nedeterministický konečný automat
\begin{itemize}
\item deterministický konečný automat (DFA) je pětice $A=(Q,\Sigma,\delta,q_0,F)$
\begin{itemize}
\item $Q$ \dots{} konečná množina stavů
\item $\Sigma$ \dots{} konečná neprázdná množina vstupních symbolů (abeceda)
\item $\delta:Q\times\Sigma\to Q$ \dots{} přechodová funkce
\item $q_0\in Q$ \dots{} počáteční stav
\item $F\subseteq Q$ \dots{} množina koncových (přijímajících) stavů
\end{itemize}
\item nedeterministický konečný automat s $\epsilon$ přechody je pětice $A=(Q,\Sigma,\delta,q_0,F)$
\begin{itemize}
\item $Q$ \dots{} konečná množina stavů
\item $\Sigma$ \dots{} konečná množina vstupních symbolů (abeceda)
\item $\delta:Q\times(\Sigma\cup\set{\epsilon})\to \mathcal P(Q)$ \dots{} přechodová funkce vracející podmnožinu $Q$
\item $q_0\in Q$ \dots{} počáteční stav
\begin{itemize}
\item alternativa: množina počátečních stavů $S_0\subseteq Q$
\end{itemize}
\item $F\subseteq Q$ \dots{} množina koncových (přijímajících) stavů
\end{itemize}
\item $\epsilon$-uzávěr \dots{} všechny stavy, kam se z daného stavu dostaneme prázdným slovem
\item podmnožinová konstrukce
\begin{itemize}
\item věta: jazyk $L$ je rozpoznatelný $\epsilon$NFA, právě když je $L$ regulární
\item pro libovolný $\epsilon$NFA $N$ zkonstruujeme DFA $D$ přijímající stejný jazyk jako $N$
\item nové stavy jsou $\epsilon$-uzavřené podmnožiny $Q_N$
\begin{itemize}
\item $Q_D\subseteq\mathcal P(Q_N)$
\item $\forall S\subseteq Q_N:\epsilon\text{closure}(S)\in Q_D$
\item v $Q_D$ může být i $\emptyset$
\end{itemize}
\item počáteční stav je $\epsilon$-uzávěr $q_0$
\begin{itemize}
\item $q_{D_0}=\epsilon\text{closure}(q_{N_0})$
\end{itemize}
\item přijímající stavy jsou všechny množiny obsahující nějaký přijímající stav
\begin{itemize}
\item $F_D=\set{S\in Q_D: S\cap F_N\neq\emptyset}$
\end{itemize}
\item přechodová funkce sjednotí a $\epsilon$-uzavře přechody z jednotlivých stavů
\begin{itemize}
\item pro $S\in Q_D,\,x\in\Sigma$ definujeme $\delta_D(S,x)=\epsilon\text{closure}\left(\bigcup_{p\in S}\delta_N(p,x)\right)$
\end{itemize}
\end{itemize}
\end{itemize}
\item Regulární výrazy
\begin{itemize}
\item regulární výrazy
\begin{itemize}
\item $\text{RegE}(\Sigma)$ \dots{} množina všech regulárních výrazů nad konečnou neprázdnou abecedou
\item $L(\alpha)$ \dots{} hodnota regulárního výrazu $\alpha$
\item regulární výraz a jeho hodnota jsou definovány induktivně
\begin{itemize}
\item základ
\begin{itemize}
\item $\epsilon$ \dots{} prázdný řetězec, $L(\epsilon)=\set{\epsilon}$
\item $\emptyset$ \dots{} prázdný výraz, $L(\emptyset)=\emptyset$
\item $a$ \dots{} znak $a\in\Sigma$, $L(a)=\set{a}$
\end{itemize}
\item indukce
\begin{itemize}
\item $\alpha+\beta$ \dots{} jako OR, $L(\alpha+\beta)=L(\alpha)\cup L(\beta)$
\item $\alpha\beta$ \dots{} $L(\alpha\beta)=L(\alpha)L(\beta)$
\item $\alpha^*$ \dots{} $L(\alpha^*)=L(\alpha)^*$
\item $(\alpha)$ \dots{} závorky nemění hodnotu, $L((\alpha))=L(\alpha)$
\end{itemize}
\end{itemize}
\item nejvyšší prioritu má iterace $^*$, pak zřetězení, pak sjednocení $+$
\item třída $\text{RegE}(\Sigma)$ je nejmenší třída uzavřená na uvedené operace
\end{itemize}
\item Kleeneho věta
\begin{itemize}
\item každý jazyk reprezentovaný konečným automatem lze zapsat jako regulární výraz
\item každý jazyk popsaný regulárním výrazem lze zapsat jako $\epsilon$NFA
\end{itemize}
\end{itemize}
\item Gramatika, jazyk generovaný gramatikou
\begin{itemize}
\item formální (generativní) gramatika je čtveřice $G=(V,T,P,S)$
\begin{itemize}
\item $V$ \dots{} konečná množina neterminálů (variables)
\item $T$ \dots{} neprázdná konečná množina terminálních symbolů (terminálů)
\item $S\in V$ \dots{} počáteční symbol
\item $P$ \dots{} konečná množina pravidel (produkcí) reprezentující rekurzivní definici jazyka
\begin{itemize}
\item každé pravidlo má tvar $\beta A\gamma\to\omega$
\begin{itemize}
\item $A\in V$
\item $\beta,\gamma,\omega\in (V\cup T)^*$
\end{itemize}
\item tj. levá strana obsahuje aspoň jeden neterminální symbol
\end{itemize}
\end{itemize}
\item jazyk $L(G)$ generovaný gramatikou $G=(V,T,P,S)$ je množina terminálních řetězců, pro které existuje derivace ze startovního symbolu $L(G)=\set{w\in T^*: S\Rightarrow_G^* w}$
\item jazyk neterminálu $A\in V$ definujeme $L(A)=\set{w\in T^*: A\Rightarrow_G^* w}$
\end{itemize}
\item Zásobníkový automat, třída jazyků přijímaných zásobníkovými automaty
\begin{itemize}
\item zásobníkový automat (PDA) je sedmice $P=(Q,\Sigma,\Gamma,\delta,q_0,Z_0,F)$
\begin{itemize}
\item $Q$ \dots{} konečná množina stavů
\item $\Sigma$ \dots{} neprázdná konečná množina vstupních symbolů
\item $\Gamma$ \dots{} neprázdná konečná zásobníková abeceda
\item $\delta$ \dots{} přechodová funkce
\begin{itemize}
\item $\delta:Q\times(\Sigma\cup\set{\epsilon})\times\Gamma\to \mathcal P_{FIN}(Q\times\Gamma^*)$
\begin{itemize}
\item kde $\mathcal P_{FIN}(A)$ je množina všech konečných podmnožin množiny $A$
\end{itemize}
\item tedy $\delta(p,a,X)\ni(q,\gamma)$
\begin{itemize}
\item $q$ je nový stav
\item $\gamma$ je řetězec zásobníkových symbolů, který nahradí $X$ na vrcholu zásobníku
\end{itemize}
\end{itemize}
\item $q_0\in Q$ \dots{} počáteční stav
\item $Z_0\in\Gamma$ \dots{} počáteční zásobníkový symbol
\item $F$ \dots{} množina přijímajících (koncových) stavů (může být nedefinovaná)
\item jak PDA funguje
\begin{itemize}
\item je nedeterministický
\item v jednom časovém kroku
\begin{itemize}
\item na vstupu přečte žádný nebo jeden symbol
\item přejde do nového stavu
\item nahradí symbol na vrchu zásobníku libovolným řetězcem (nejlevější znak bude na vrchu)
\end{itemize}
\end{itemize}
\item přechodový diagram
\begin{itemize}
\item jako pro konečný automat
\item hrany popisujeme ve tvaru: \textit{vstupní znak}, \textit{zásobníkový znak} $\to$ \textit{push řetězec}
\end{itemize}
\end{itemize}
\item jazyk přijímaný koncovým stavem je $L(P_{da})=\set{w\in\Sigma^*:(\exists q\in F)(\exists\alpha\in\Gamma^*)((q_0,w,Z_0)\vdash^*_{P_{da}}(q,\epsilon,\alpha))}$
\item jazyk přijímaný prázdným zásobníkem je $N(P_{da})=\set{w\in\Sigma^*:(\exists q\in Q)((q_0,w,Z_0)\vdash^*_{P_{da}}(q,\epsilon,\epsilon))}$
\begin{itemize}
\item množina $F$ je nerelevantní, takže ji lze vynechat - pak je PDA šestice
\end{itemize}
\item lemma: od přijímajícího stavu k prázdnému zásobníku
\begin{itemize}
\item na dno zásobníku přidáme špunt
\item z přijímajících stavů vedeme hranu do vypouštěcího stavu, kde zásobník vyprázdníme
\end{itemize}
\item lemma: od prázdného zásobníku ke koncovému stavu
\begin{itemize}
\item na dno zásobníku přidáme špunt
\item z každého stavu uděláme přechod takový, že pokud je vidět špunt, přemístíme se do koncového stavu
\end{itemize}
\item věta: následující tvrzení jsou ekvivalentní
\begin{itemize}
\item jazyk $L$ je bezkontextový, tj. generovaný bezkontextovou gramatikou
\item jazyk $L$ je přijímaný nějakým zásobníkovým automatem koncovým stavem
\item jazyk $L$ je přijímaný nějakým zásobníkovým automatem prázdným zásobníkem
\end{itemize}
\item konstrukce PDA z CFG
\begin{itemize}
\item mějme gramatiku $G=(V,T,P,S)$
\item konstruujeme PDA $P=(\set{q},T,V\cup T,\delta,q,S)$
\begin{itemize}
\item $\forall A\in V:\delta(q,\epsilon,A)=\set{(q,\beta)\mid (A\to \beta)\in P}$
\item $\forall a\in T:\delta(q,a,a)=\set{(q,\epsilon)}$
\item $P$ přijímá prázdným zásobníkem
\end{itemize}
\item idea
\begin{itemize}
\item stavy nás nezajímají - stačí nám jeden
\item na zásobníku nedeterministicky generujeme všechny možné posloupnosti terminálů (přičemž ale postupně zpracováváme vstup)
\end{itemize}
\end{itemize}
\end{itemize}
\item Pumping lemma pro bezkontextové jazyky
\begin{itemize}
\item lemma o vkládání (pumping) pro bezkontextové jazyky
\begin{itemize}
\item nechť $L$ je bezkontextový jazyk
\item $(\exists n)(\forall w\in L):|w|\geq n\implies w=u_1u_2u_3u_4u_5$
\begin{itemize}
\item $u_2u_4\neq\epsilon$
\item $|u_2u_3u_4|\leq n$
\item $\forall k\geq 0:u_1u_2^ku_3u_4^ku_5\in L$
\end{itemize}
\end{itemize}
\item idea důkazu
\begin{itemize}
\item vezmeme derivační strom pro $w$
\item nejdeme nejdelší cestu, na ní dva stejné neterminály
\item tyto neterminály určí dva podstromy, které definují rozklad slova
\item větší podstrom můžeme posunout ($k\gt 1$) nebo nahradit menším podstromem ($k=0$)
\end{itemize}
\item příklad: $L_{012}=\set{0^i1^i2^i\mid i\geq 1}$ není bezkontextový
\begin{itemize}
\item důkaz sporem - předpokládejme bezkontextovost
\item vezmeme $n$ z pumping lemmatu
\item zvolíme $w=0^n1^n2^n$
\item délka pumpovacího slova $u_2u_3u_4$ musí být nejvýše $n$, tedy vždy můžeme pumpovat maximálně dva různé symboly
\item $u_2u_4\neq\epsilon\implies$ iterací se slovo změní
\item tím se poruší rovnost symbolů
\item tedy ať už zvolíme libovolné dělení $u_1u_2u_3u_4u_5$, nové slovo nebude patřit do $L_{012}$, což je spor
\end{itemize}
\end{itemize}
\item Chomského normální tvar
\begin{itemize}
\item o bezkontextové gramatice $G=(V,T,P,S)$ bez zbytečných symbolů, kde jsou všechna pravidla ve tvaru $A\to BC$ nebo $A\to a$, kde $A,B,C\in V,\, a\in T$, říkáme, že je v Chomského normálním tvaru (ChNF)
\item algoritmus převodu
\begin{itemize}
\item eliminujeme $\epsilon$-pravidla
\begin{itemize}
\item označíme nulovatelné symboly
\item přidáme verzi pravidel bez nulovatelných symbolů
\item odstraníme $\epsilon$-pravidla
\end{itemize}
\item eliminujeme jednotková pravidla (tím nepřidáme $\epsilon$-pravidla)
\begin{itemize}
\item najdeme jednotkové páry
\item přidáme odpovídající pravidla
\item odstraníme jednotková pravidla
\end{itemize}
\item eliminujeme zbytečné symboly (tím nepřidáme žádná pravidla)
\begin{itemize}
\item nejdříve eliminujeme negenerující, pak nedosažitelné
\end{itemize}
\item pravé strany délky aspoň 2 předěláme na samé neterminály
\item pravé strany s aspoň třemi neterminály rozdělíme na více pravidel
\end{itemize}
\end{itemize}
\item Turingův stroj
\begin{itemize}
\item turingův stroj (TM) je sedmice $M=(Q,\Sigma,\Gamma,\delta,q_0,B,F)$
\begin{itemize}
\item $Q$ \dots{} konečná množina stavů
\item $\Sigma$ \dots{} konečná neprázdná množina vstupních symbolů
\item $\Gamma$ \dots{} konečná množina všech symbolů pro pásku
\begin{itemize}
\item vždy $\Gamma\supseteq\Sigma,\,Q\cap\Gamma=\emptyset$
\end{itemize}
\item $\delta$ \dots{} (částečná) přechodová funkce
\begin{itemize}
\item zobrazení $(Q-F)\times\Gamma\to Q\times\Gamma\times\set{L,R}$
\item $\delta(q,X)=(p,Y,D)$
\begin{itemize}
\item $q\in (Q-F)$ \dots{} aktuální stav
\item $X\in \Gamma$ \dots{} aktuální symbol na pásce
\item $p\in Q$ \dots{} nový stav
\item $Y\in\Gamma$ \dots{} symbol pro zapsání do aktuální buňky, přepíše aktuální obsah
\item $D\in\set{L,R}$ \dots{} směr pohybu hlavy (doleva, doprava)
\end{itemize}
\end{itemize}
\item $q_0\in Q$ \dots{} počáteční stav
\item $B\in\Gamma-\Sigma$ \dots{} blank = symbol pro prázdné buňky, na začátku všude kromě konečného počtu buněk se vstupem
\item $F\subseteq Q$ \dots{} množina koncových neboli přijímajících stavů
\item poznámka: někdy se nerozlišuje $\Gamma$ a $\Sigma$ a neuvádí se $B$, pak je TM pětice
\end{itemize}
\item turingův stroj $M=(Q,\Sigma,\Gamma,\delta,q_0,B,F)$ přijímá jazyk $L(M)=\set{w\in\Sigma^*:(\exists p\in F)(\exists \alpha,\beta\in\Gamma^*)(q_0w\vdash_M^*\alpha p\beta)}$
\begin{itemize}
\item tj. množinu slov, po jejichž přečtení se dostane do koncového stavu
\item pásku nemusí uklízet (v naší definici)
\end{itemize}
\item jazyk nazveme rekurzivně spočetným, pokud je přijímán nějakým Turingovým strojem
\item říkáme, že TM $M$ rozhoduje jazyk $L$, pokud $L=L(M)$ a pro každé $w\in\Sigma^*$ stroj nad $w$ zastaví
\begin{itemize}
\item jazyky rozhodnutelné TM nazýváme rekurzivní jazyky
\item rekurzivní jazyky jsou podmnožinou rekurzivně spočetných jazyků
\end{itemize}
\end{itemize}
\item Algoritmicky nerozhodnutelné problémy
\begin{itemize}
\item rozhodnutelný problém
\begin{itemize}
\item problém \dots{} množina otázek (nebo spíše vstupů) kódovatelná řetězci nad abecedou $\Sigma$ s odpověďmi $\in\set{\text{ano},\text{ne}}$
\begin{itemize}
\item pokud problém definuji jako množinu, jde o otázku, zda vstup kóduje prvek dané množiny (např. polynom s celočíselným kořenem)
\end{itemize}
\item problém je (algoritmicky) rozhodnutelný, pokud existuje TM takový, že pro každý vstup $w\in P$ TM zastaví a navíc přijme, právě když $P(w)=\text{ano}$ (tj. pro $P(w)=\text{ne}$ zastaví v nepřijímajícím stavu)
\item nerozhodnutelný problém \dots{} není rozhodnutelný
\item existuje analogie mezi rozhodnutelným problémem a rekurzivním jazykem
\end{itemize}
\item redukce problému
\begin{itemize}
\item definice: redukcí problému $P_1$ na $P_2$ nazýváme algoritmus $R$, který pro každou instanci $w\in P_1$ zastaví a vydá $R(w)\in P_2$, přičemž\dots{}
\begin{itemize}
\item $P_1(w)=\text{ano}\iff P_2(R(w))=\text{ano}$
\item tedy i $P_1(w)=\text{ne}\iff P_2(R(w))=\text{ne}$
\end{itemize}
\item věta: existuje-li redukce problému $P_1$ na $P_2$, pak\dots{}
\begin{itemize}
\item pokud $P_1$ je nerozhodnutelný, pak je nerozhodnutelný i $P_2$
\item pokud $P_1$ není rekurzivně spočetný, pak není rekurzivně spočetný ani $P_2$
\end{itemize}
\end{itemize}
\item problém zastavení není rozhodnutelný
\begin{itemize}
\item definice: instancí problému zastavení je dvojice řetězců $M,w\in\set{0,1}^*$
\item definice: problém zastavení je najít algoritmus $\text{Halt}(M,w)$, který vydá 1, právě když stroj $M$ zastaví na vstupu $w$, jinak vydá 0
\item věta: problém zastavení není rozhodnutelný
\item idea důkazu: redukujeme $L_d$ na $\text{Halt}$
\end{itemize}
\item diagonální jazyk
\begin{itemize}
\item diagonální jazyk $L_d$ je definovaný jako jazyk slov $w\in\set{0,1}^*$ takových, že TM reprezentovaný jako $w$ nepřijímá slovo $w$
\item věta: $L_d$ není rekurzivně spočetný jazyk, tj. neexistuje TM přijímající $L_d$
\item idea důkazu
\begin{itemize}
\item kdyby existoval TM přijímající $L_d$, spuštění takového stroje na vlastním kódu by vedlo k paradoxu
\end{itemize}
\end{itemize}
\item univerzální jazyk $L_u$ definujeme jako množinu binárních řetězců, které kódují pár $(M,w)$, kde $M$ je TM a $w\in L(M)$
\item TM přijímající $L_u$ se nazývá \textit{univerzální Turingův stroj}
\begin{itemize}
\item existuje
\end{itemize}
\item problém univerzálního jazyka není rozhodnutelný
\begin{itemize}
\item věta: $L_u$ je rekurzivně spočetný, ale není rekurzivní
\item mohli bychom zkonstruovat TM přijímající $L_d$
\end{itemize}
\item problém prázdnosti jazyka daného TM není rozhodnutelný
\end{itemize}
\item Chomského hierarchie
\begin{itemize}
\item gramatiky typu 0 (rekurzivně spočetné jazyky $\mathcal L_0$)
\begin{itemize}
\item pravidla v obecné formě $\alpha\to\omega$
\begin{itemize}
\item $\alpha,\omega\in(V\cup T)^*$
\item $\alpha$ obsahuje neterminál
\end{itemize}
\item rekurzivně spočetný jazyk je rozpoznatelný (nějakým) Turingovým strojem
\end{itemize}
\item gramatiky typu 1 (kontextové gramatiky, jazyky $\mathcal L_1$)
\begin{itemize}
\item pouze pravidla ve tvaru $\gamma A\beta\to\gamma\omega\beta$
\begin{itemize}
\item $A\in V$
\item $\gamma,\beta\in(V\cup T)^*$
\item $\omega\in (V\cup T)^+$
\item jedinou výjimkou je pravidlo $S\to\epsilon$, pak se ale $S$ nevyskytuje na pravé straně žádného pravidla
\end{itemize}
\item kontextový jazyk je rozpoznatelný lineárně omezeným automatem (LBA, lineárně omezený Turingův stroj)
\end{itemize}
\item gramatiky typu 2 (bezkontextové gramatiky, jazyky $\mathcal L_2$)
\begin{itemize}
\item pouze pravidla ve tvaru $A\to\omega$
\begin{itemize}
\item $A\in V$
\item $\omega\in(V\cup T)^*$
\end{itemize}
\item bezkontextový jazyk je rozpoznatelný nedeterministickým zásobníkovým automatem
\end{itemize}
\item gramatiky typu 3 (regulární gramatiky / pravé lineární gramatiky, regulární jazyky $\mathcal L_3$)
\begin{itemize}
\item pouze pravidla ve tvaru $A\to\omega B$ nebo $A\to\omega$
\begin{itemize}
\item $A,B\in V$
\item $\omega\in T^*$
\end{itemize}
\item idea
\begin{itemize}
\item každý neterminál odpovídá stavu konečného automatu
\item pravidla odpovídají přechodové funkci
\end{itemize}
\item regulární jazyk je rozpoznatelná konečným automatem
\end{itemize}
\end{itemize}
\item Schopnost zařazení konkrétního jazyka do Chomského hierarchie (zpravidla sestrojení odpovídajícího automatu či gramatiky)
\begin{itemize}
\item obvykle chceme jazyk zařadit co nejpřesněji - např. říct, že není regulární (dokázat pomocí iteračního lemmatu) a že je bezkontextový (sestavit CFG)
\item je regulární: sestrojíme konečný automat
\item není regulární: použijeme iterační lemma nebo Myhill-Nerodovu větu
\item je bezkontextový: sestrojíme bezkontextovou gramatiku
\item není bezkontextový: použijeme lemma o vkládání
\item kontextovostí se obvykle nezabýváme
\begin{itemize}
\item lemma: jazyk $L=\set{a^ib^jc^k\mid1\leq i\leq j\leq k}$ je kontextový jazyk, není bezkontextový
\end{itemize}
\item je rekurzivně spočetný: popíšeme algoritmus
\begin{itemize}
\item algoritmus pro $0^n1^n$
\begin{itemize}
\item jezdíme tam a zpátky po pásce, přičemž $0$ přepíšeme na $X$, znak $1$ přepíšeme na $Y$
\item přijmeme, pokud nám jedničky a nuly dojdou ve stejnou chvíli (páska je popsaná samými $X$ a $Y$)
\end{itemize}
\end{itemize}
\item není rekurzivně spočetný (nebo případně jenom rekurzivní): dokážeme pomocí $L_d$
\item uzavřenost operací
\begin{itemize}
\item regulární jazyky jsou uzavřené na množinové (sjednocení, průnik, rozdíl, doplněk) i řetězcové operace ($L{.}M,L^*,L^+,L^R,M\setminus L,L/M$), na homomorfismus a inverzní homomorfismus
\begin{itemize}
\item doplněk: obrátíme \quotedblbase{}koncovost\textquotedblleft{} stavů
\item průnik, sjednocení, rozdíl
\begin{itemize}
\item zkonstruujeme součinový automat $(Q_1\times Q_2,\Sigma,\delta,(q_{0_1},q_{0_2}),F)$
\begin{itemize}
\item kde $\delta((p,q),x)=(\delta_1(p,x),\delta_2(q,x))$
\end{itemize}
\item průnik \dots{} $F=F_1\times F_2$
\item sjednocení \dots{} $F=(F_1\times Q_2)\cup(Q_1\times F_2)$
\item rozdíl \dots{} $F=F_1\times (Q_2-F_2)$
\end{itemize}
\end{itemize}
\item bezkontextové jazyky jsou uzavřené na sjednocení, konkatenaci, iteraci, reverzi, naopak neuzavřené na průnik (ale jsou uzavřené na průnik s regulárním jazykem)
\item uzavřenost můžeme použít k důkazu regularity nebo bezkontextovosti
\item ale můžeme ji použít i k důkazu neregularity jazyka $L$ - kdybychom aplikací uzavřených operací na jazyk $L$ a regulární jazyky dostali jazyk $L'$, který zjevně není regulární (např. $0^i1^i$)
\end{itemize}
\end{itemize}
\item Příklad kontextového jazyka
\begin{itemize}
\item lemma: jazyk $L=\set{a^ib^jc^k\mid1\leq i\leq j\leq k}$ je kontextový jazyk, není bezkontextový
\item pravidla
\begin{itemize}
\item $S\to aSBC|aBC$ \dots{} generování symbolů $a$
\item $B\to BBC$ \dots{} množení symbolů $B$
\item $C\to CC$ \dots{} množení symbolů $C$
\item $CB\to BC$ \dots{} uspořádání symbolů $B$ a $C$
\begin{itemize}
\item pravidlo není kontextové, je potřeba ho upravit nadtrháváním
\item tzn. písmenka měníme po jednom (střídavě), např. takto: $CB\to\bar CB\to \bar C\bar B\to B\bar B\to BC$
\end{itemize}
\item $aB\to ab$ \dots{} začátek přepisu $B$ na $b$
\item $bB\to bb$ \dots{} pokračování přepisu $B$ na $b$
\item $bC\to bc$ \dots{} začátek přepisu $C$ na $c$
\item $cC\to cc$ \dots{} pokračování přepisu $C$ na $c$
\end{itemize}
\item je důležité, že tam chybí $cB\to cb$
\end{itemize}
\end{itemize}
\subsection{2. Algoritmy a datové struktury}
\begin{itemize}
\item Časová a prostorová složitost algoritmu
\begin{itemize}
\item pro různé úlohy používáme různé parametrizace vstupu (tedy velikost vstupu může být trochu něco jiného - někdy počet čísel, někdy jejich hodnota\dots{})
\item doba běhu pro vstup x
\begin{itemize}
\item $t(x):=$ součet cen provedených instrukcí (může být $\infty$)
\end{itemize}
\item časová složitost výpočtu
\begin{itemize}
\item $T(n):= \text{max}\lbrace t(x)\mid x \text{ vstup velikosti }n\rbrace$
\end{itemize}
\item prostor běhu pro vstup x
\begin{itemize}
\item $s(x):=$ max. adresa - min. adresa + 1
\item máme na mysli maximální a minimální adresy navštívené během výpočtu
\end{itemize}
\item prostorová složitost výpočtu
\begin{itemize}
\item $S(n):= \text{max}\lbrace s(x)\mid x \text{ vstup velikosti }n\rbrace$
\end{itemize}
\item takhle jsme definovali složitosti v nejhorším případě, někdy se může hodit i průměrný případ
\end{itemize}
\item Měření velikosti dat
\begin{itemize}
\item spotřebovanou paměť definujeme jako rozdíl mezi nejvyšším a nejnižším použitým indexem paměti
\item paměťová složitost pak odpovídá maximu ze spotřeby paměti přes všechny vstupy dané velikosti
\item do časové složitosti nepočítáme načtení vstupu
\item podobně do prostorové složitosti nepočítáme vstup, do nějž se nezapisuje, ani výstup, pokud se jenom zapíše a pak už se nečte (takhle se někdy definuje pracovní prostor výpočtu)
\item je potřeba omezit kapacitu paměťové buňky
\begin{itemize}
\item jednotková cena instrukce - omezíme šířku slova na W bitů
\item logaritmická cena instrukce - cena odpovídá počtu bitů operandů včetně adres (je přesný, ale nepohodlný na přemýšlení o algoritmech)
\item relativní logaritmická cena instrukce - u polynomiálně velkých čísel je cena jednotková, u obrovských čísel se cena zvyšuje (tohle budeme reálně používat - i když se k tomu budeme obvykle chovat jako k jednotkové ceně instrukce)
\end{itemize}
\end{itemize}
\item Složitost v nejlepším, nejhorším a průměrném případě
\begin{itemize}
\item nejčastěji se používá složitost v nejhorším případě
\begin{itemize}
\item složitost pro nejhorší možný vstup
\end{itemize}
\item složitost v průměrném případě dobře charakterizuje kvalitu algoritmu, ale je obtížné ji odvodit
\begin{itemize}
\item průměrujeme přes všechny možné vstupy
\end{itemize}
\end{itemize}
\item Asymptotická notace
\begin{itemize}
\item $f\in O(g)\equiv\exists c\,\forall^*n\;f(n)\leq c\cdot g(n)$
\begin{itemize}
\item $f,g:\mathbb N\to\mathbb R$
\item $\forall^*n$ \dots{} pro všechna $n$ až na konečně mnoho výjimek (existuje $n_0$ takové, že pro všechna větší $n$ už to platí)
\end{itemize}
\item $f\in \Omega(g)\equiv\exists c\,\forall^*n\;f(n)\geq c\cdot g(n)$
\item $\Theta(g):= O(g)\cap \Omega(g)$
\end{itemize}
\item Třídy P a NP
\begin{itemize}
\item problém $L\in \text P\equiv\exists A$ algoritmus $\exists p$ polynom takový, že $\forall x$ vstup platí, že $A(x)$ doběhne do $p(|x|)$ kroků $\land\;A(x)=L(x)$
\item problém $L\in \text{NP}\equiv\exists V\in \text P$ (verifikátor) $\exists g$ polynom (omezení délky certifikátů) $\forall x:L(x)=1\iff$ $\exists y$ certifikát $|y|\leq g(|x|)$ (krátký) $\land\;V(x,y)=1$ (schválený)
\item $\text P\subseteq\text{NP}$ (rovnost se neví)
\end{itemize}
\item Převoditelnost problémů, NP-těžkost a NP-úplnost
\begin{itemize}
\item rozhodovací problém $\equiv$ funkce $f:\set{0,1}^*\to\set{0,1}$
\item převod
\begin{itemize}
\item mějme rozhodovací problémy $A,B$
\item problém $A$ je převoditelný na problém $B$ právě tehdy, když existuje funkce $f:\set{0,1}^*\to\set{0,1}^*$ taková, že $\forall\alpha\in\set{0,1}^*:A(\alpha)=B(f(\alpha))$ a $f$ lze spočítat v čase polynomiálním vzhledem k $|\alpha|$
\item značíme $A\to B$ nebo $A\leq_P B$
\item funkci $f$ říkáme převod (případně redukce)
\end{itemize}
\item vlastnosti převoditelnosti
\begin{itemize}
\item je reflexivní $(A\to A)$ \dots{} $f$ je identita
\item je tranzitivní $(A\to B\land B\to C\implies A\to C)$ \dots{} když $f$ převádí $A$ na $B$ a $g$ převádí $B$ na $C$, tak $f\circ g$ převádí $A$ na $C$ (přičemž složení polynomiálně vyčíslitelných funkcí je polynomiálně vyčíslitelná funkce)
\item není antisymetrická - např. problémy \quotedblbase{}vstup má sudou délku\textquotedblleft{} a \quotedblbase{}vstup má lichou délku\textquotedblleft{} lze mezi sebou převádět oběma směry
\item existují navzájem nepřevoditelné problémy - např. mezi problémy \quotedblbase{}na každý vstup odpověz 0\textquotedblleft{} a \quotedblbase{}na každý vstup odpověz 1\textquotedblleft{} nemůže existovat převod
\item převoditelnost je částečné kvaziuspořádání na množině všech problémů
\end{itemize}
\item definice: NP-těžké a NP-úplné problémy
\begin{itemize}
\item $L$ je NP-těžký $\equiv\forall K\in\text{NP}:K\to L$
\item $L$ je NP-úplný $\equiv L$ je NP-těžký $\land\;L\in\text{NP}$
\end{itemize}
\item věta: pokud $A\to B$ a $B\in \text P$, pak $A\in \text P$
\item věta: pokud $A\to B$, $B\in \text{NP}$ a $A$ je NP-úplný, pak $B$ je NP-úplný
\end{itemize}
\item Příklady NP-úplných problémů a převodů mezi nimi
\begin{itemize}
\item logické: (CNF) SAT, 3-SAT, 3,3-SAT, obvodový SAT
\item grafové: nezávislá množina, klika, $k$-obarvitelnost (pro $k\geq 3$), hamiltonovská cesta/kružnice, 3D-párování
\item číselné: součet podmnožiny, batoh, 2 loupežníci, nulajedničkové řešení soustavy lineárních rovnic (v $\mathbb Z$)
\item popisy problémů
\begin{itemize}
\item klika: existuje úplný podgraf grafu $G$ na alespoň $k$ vrcholech?
\item nezávislá množina: existuje nezávislá množina vrcholů grafu $G$ velikosti aspoň $k$?
\begin{itemize}
\item množina vrcholů grafu je nezávislá $\equiv$ žádné dva vrcholy ležící v této množině nejsou spojeny hranou
\end{itemize}
\item SAT: existuje dosazení 0 a 1 za proměnné tak, aby $\psi(\dots)=1$?
\begin{itemize}
\item kde $\psi$ je v CNF
\end{itemize}
\item 3-SAT: jako SAT, akorát každá klauzule formule $\psi$ obsahuje nejvýše tři literály
\item 3,3-SAT: jako 3-SAT, akorát se každá proměnná vyskytuje v~maximálně třech literálech
\item 3D-párování
\begin{itemize}
\item vstup: tři množiny $A,B,C$ a množina kompatibilních trojic $T\subseteq A\times B\times C$
\item výstup: existuje perfektní podmnožina trojic? tzn. existuje taková podmnožina, v níž se každý prvek množin $A,B,C$ účastní právě jedné trojice
\end{itemize}
\end{itemize}
\item převod klika $\leftrightarrow$ nezávislá množina
\begin{itemize}
\item pokud v grafu prohodíme hrany a nehrany, stane se z každé kliky nezávislá množina a naopak
\item převodní funkce zneguje hrany
\end{itemize}
\item při převodu problémů typu SAT vyrábíme ekvisplnitelnou formuli
\item SAT $\to$ 3-SAT
\begin{itemize}
\item $(\alpha\lor\beta)\to(\alpha\lor\zeta)\land(\beta\lor\neg\zeta)$
\item převod funguje i naopak (viz rezoluce)
\item jak rozštípnout dlouhou klauzuli délky $\ell$?
\begin{itemize}
\item $\alpha$ nechť má délku 2
\item $\beta$ nechť má délku $\ell-2$
\item po přidání $\zeta$ dostanu konjunkci klauzulí délky 3 a $\ell-1$
\item klauzuli délky $\ell-1$ štípu dál (pokud je moc dlouhá)
\end{itemize}
\item počet štípnutí je shora omezen délkou formule
\item v polynomiálním čase postupně rozštípeme všechny dlouhé klauzule při zachování splnitelnosti
\end{itemize}
\item 3-SAT $\to$ 3,3-SAT
\begin{itemize}
\item nechť $x$ je proměnná s $k\gt 3$ výskyty
\item pro každý výskyt si pořídíme nové proměnné $x_1,\dots,x_k$
\item ekvivalenci všech $x_i$ zajistíme řetězcem implikací
\begin{itemize}
\item $x_1\implies x_2$
\item $x_2\implies x_3$
\item $\quad\vdots$
\item $x_k\implies x_1$
\end{itemize}
\item každá proměnná $x_i$ se tudíž vyskytne třikrát
\end{itemize}
\item 3,3-SAT* \dots{} navíc každý literál max. 2$\times$
\begin{itemize}
\item použijeme předchozí algoritmus pro všechny proměnné s $k \geq 3$ výskyty
\end{itemize}
\item 3-SAT $\to$ nezávislá množina
\begin{itemize}
\item pro každou z $k$ klauzulí formule vytvoříme trojúhelník a jeho vrcholům přiřadíme literály klauzule
\begin{itemize}
\item pokud klauzule obsahuje méně než tři literály, tak přebývající vrcholy smažeme
\end{itemize}
\item spojíme hranami všechny dvojice konfliktních literálů
\item v takovém grafu existuje nezávislá množina velikosti alespoň $k$ právě tehdy, když je formule splnitelná
\begin{itemize}
\item máme-li splňující ohodnocení, můžeme z každé klauzule vybrat jeden pravdivý literál - ty umístíme do nezávislé množiny
\item máme-li nezávislou množinu velikosti $k$, vybereme literály odpovídající těmto vrcholům a podle nich nastavíme odpovídající proměnné
\begin{itemize}
\item v nezávislé množině velikosti $k$ nemůžou být dva vrcholy ze stejného trojúhelníku - tedy z každého trojúhelníku (klauzule) bude v množině právě jeden vrchol
\item v nezávislé množině nemohou být dva vrcholy s opačnými literály, protože takové vrcholy jsou spojeny hranou
\end{itemize}
\end{itemize}
\end{itemize}
\item nezávislá množina $\to$ 3-SAT
\begin{itemize}
\item každému vrcholu odpovídá jedna proměnná (pokud je proměnná ohodnocena jedničkou, vrchol náleží do nezávislé množiny)
\item pro každou hranu přidáme klauzuli $(\neg x_i\lor\neg x_j)$, která zajistí, že obě proměnné nebudou splněny zároveň
\item chceme nezávislou množinu velikosti $k$, takže musíme vrcholy v množině spočítat
\begin{itemize}
\item pořídíme si tabulku (pomocí jiných proměnných $y_{ij}$)
\item vynutíme, aby v každém sloupci byla maximálně jedna jednička a v každém řádku právě jedna jednička
\end{itemize}
\item propojíme proměnné $x_j$ a $y_{ij}$ pomocí implikace $y_{ij}\to x_j$
\end{itemize}
\end{itemize}
\item Metoda rozděl a panuj: princip rekurzivního dělení problému na podproblémy
\begin{itemize}
\item problém rekurzivně dělíme na menší a menší podproblémy, dokud nedojdeme k problému konstantní velikosti, který umíme vyřešit triviálně
\item problémy v jedné úrovni dělení na sobě musí být nezávislé
\end{itemize}
\item Metoda rozděl a panuj: výpočet složitosti pomocí rekurentních rovnic
\begin{itemize}
\item vyjádříme složitost algoritmu (tedy $T(n)$) pomocí složitosti podproblému, tím dostaneme rekurzivní rovnici
\item do rovnice můžeme postupně dosazovat a vypozorovat nějaký obecný vzorec - chceme dojít až k $T(1)$
\item nebo můžeme rovnou použít Master theorem
\item případně ze složitosti $i$-té úrovně můžeme odvodit složitost celého algoritmu (součtem přes všechny úrovně)
\end{itemize}
\item Master theorem (kuchařková věta) - bez důkazu
\begin{itemize}
\item uvažujeme rekurzivní algoritmus, který vstup rozloží na $a$ podproblémů velikosti $n/b$ a z jejich výsledků složí celkovou odpověď v čase $\Theta(n^c)$
\item věta: rekurentní rovnice $T(n)=a\cdot T(n/b)+\Theta(n^c),\;T(1)=1$ má pro konstanty $a\in\set{1,2,\dots},\,b\in(1,\infty),\,c\in[0,\infty)$ řešení\dots{}
\begin{itemize}
\item $T(n)=\Theta(n^c\log n)$, pokud $a/b^c=1$
\item $T(n)=\Theta(n^c)$, pokud $a/b^c\lt 1$
\item $T(n)=\Theta(n^{\log_b a})$, pokud $a/b^c\gt 1$
\end{itemize}
\end{itemize}
\item Metoda rozděl a panuj: aplikace (Mergesort, násobení dlouhých čísel)
\begin{itemize}
\item Mergesort
\begin{itemize}
\item na vstupu posloupnost
\item pokud je délka posloupnosti menší rovna jedné, mergesort vrátí vstup
\item jinak vrátí merge(mergesort(první polovina vstupu), mergesort(druhá polovina vstupu)), kde merge slévá dvě poloviční setříděné posloupnosti v lineárním čase
\item čas $\Theta(n\log n)$
\begin{itemize}
\item v každém z $\log n$ pater (lineárně) sléváme kousky posloupnosti
\item v každém patře se kousky sečtou na $n$
\end{itemize}
\end{itemize}
\item násobení $n$-ciferných čísel v čase $O(n^{\log_23})$
\begin{itemize}
\item Karacubovo násobení
\item násobíme dvě $n$-ciferná čísla $x,y$ v soustavě o základu $z$
\item cifry obou čísel rozdělíme na dvě poloviny
\begin{itemize}
\item $x=a\cdot z^{n/2}+b$
\item $y=c\cdot z^{n/2}+d$
\end{itemize}
\item pak $x\cdot y=(a\cdot z^{n/2}+b)(c\cdot z^{n/2}+d)=ac\cdot z^{n}+(ad+bc)\cdot z^{n/2}+bd$
\item $(ad+bc)$ lze vyjádřit jako $(a+b)(c+d)-ac-bd$
\item takže stačí 3 násobení
\item $T(n)=3\cdot T(n/2)+cn$
\item časová složitost podproblému na vrstvě a počet podproblémů
\begin{itemize}
\item 1\textbackslash{}. vrstva $\quad n\quad1$
\item 2\textbackslash{}. vrstva $\quad\frac n2\quad3$
\item $i$-tá vrstva $\quad \frac {n}{2^i}\quad 3^i$
\item poslední vrstva $\quad 1\quad 3^{\log_2 n}$
\end{itemize}
\item $T(n)=\sum_{i=0}^{\log_2 n}n\cdot(3/2)^i\in\Theta(n\cdot (3/2)^{\log_2n})$, což lze převést na $\Theta(n^{\log_23})$
\item podle Master theorem $a=3,\,b=2,\,c=1$
\end{itemize}
\end{itemize}
\item Definice (binárního) vyhledávacího stromu
\begin{itemize}
\item BVS je binární strom, jehož každému vrcholu $v$ přiřadíme unikátní klíč $k(v)$ z univerza. Přitom musí pro každý vrchol $v$ platit:
\begin{itemize}
\item kdykoliv $a\in L(v)$, pak $k(a)\lt k(v)$
\item kdykoliv $b\in R(v)$, pak $k(b)\gt k(v)$
\end{itemize}
\end{itemize}
\item Operace s nevyvažovanými binárními vyhledávacími stromy
\begin{itemize}
\item Find: \quotedblbase{}binární vyhledávání\textquotedblleft{} (porovnávám s vrcholem - podle toho se zanořím do jednoho z podstromů nebo vrátím hodnotu vrcholu, případně $\emptyset$, pokud se nemám kam zanořit)
\item Insert: vložím vrchol tam, kde bych ho našel, kdyby ve stromu byl (pokud tam je, tak nic nedělám)
\item Delete: vrchol najdeme a smažeme; pokud má jednoho syna, tak ho připojíme pod otce smazaného vrcholu; pokud má dva syny, nejdříve nalezneme nejlevější vrchol v pravém podstromu a s tím ho prohodíme (fungovalo by to i symetricky)
\end{itemize}
\item AVL stromy (definice)
\begin{itemize}
\item AVL strom = hloubkově vyvážený strom
\item pro každý jeho vrchol platí $|h(L(v))-h(R(v))|\leq1$
\item tedy hloubka levého a pravého podstromu se liší nejvýše o jedna
\item věta: AVL strom na $n$ vrcholech má hloubku $\Theta(\log n)$
\end{itemize}
\item Primitivní třídicí algoritmy (Bubblesort, Insertsort)
\begin{itemize}
\item BubbleSort = výměna dvou prvků
\begin{itemize}
\item procházíme pole zleva doprava
\item bereme dvojice prvků, které jsou v poli vedle sebe - pokud jsou ve špatném pořadí, tak je prohodíme
\item pokud jsme během jednoho průchodu provedli aspoň jedno prohození, musíme pole projít znova (jinak algoritmus končí)
\end{itemize}
\item InsertSort = zařazení prvku mezi již setříděné
\begin{itemize}
\item rozdělíme pole na setříděnou a nesetříděnou část (na začátku je nesetříděná část prázdná)
\item postupně bereme prvky z nesetříděné části a zařazujeme je na správné místo v setříděné části (najdeme vhodné místo, setříděné prvky vpravo posuneme napravo, do vzniklého místa vložíme aktuální prvek)
\end{itemize}
\end{itemize}
\item Quicksort
\begin{itemize}
\item určíme pivota, rozdělíme na prvky menší rovny a větší rovny
\item algoritmus $\mathrm{QuickSort}(X)$
\begin{itemize}
\item pokud $n\leq 1$: vrátíme vstup
\item zvolíme pivota $p$
\item rozdělíme $x_1,\dots,x_n$ podle $p$ na $L,S,P$
\item $L\leftarrow \text{QuickSort}(L)$
\item $P\leftarrow \text{QuickSort}(P)$
\item vrátíme $L,S,P$
\end{itemize}
\item věta: QuickSort s náhodnou volbou pivota má časovou složitost se střední hodnotou $\Theta(n\log n)$
\item částečné zdůvodnění průměrné složitosti
\begin{itemize}
\item s pravděpodobností $\frac12$ bude náhodně vybraný pivot v prostředních dvou čtvrtinách setříděné posloupnosti (\quotedblbase{}skoromedián\textquotedblleft{})
\item střední hodnota počtu pokusů, než zvolíme skoromedián, je podle lemmatu o džbánu rovna $2$ (to se schová do konstanty)
\item když jsme zvolili skoromedián, tak se problém zmenšil na $\frac34$
\end{itemize}
\item složitost v nejhorším případě je $\Theta(n^2)$
\end{itemize}
\item Dolní odhad složitosti porovnávacích třídicích algoritmů
\begin{itemize}
\item uvažujeme strom algoritmu - jeho listy odpovídají možným výsledkům
\item větvení v rozhodovacím stromě budou odpovídat jednotlivým porovnáním
\item listů musí být $n!$
\item hloubka stromu musí být aspoň $\log_2n!\geq \log_2 n^{\frac n2}=\frac n2 \log n$
\item každý každý porovnávací třídicí algoritmus musí mít složitost $\in\Omega(n\log n)$
\end{itemize}
\item Prohledávání grafu do šířky a do hloubky
\begin{itemize}
\item prohledávání do šířky (BFS)
\begin{itemize}
\item zpracování určitého vrcholu = odeberu ho z fronty, podívám se na sousední vrcholy a pokud jsou nenavštívené, tak je označím jako otevřené a přidám je do fronty ke zpracování, zpracovávaný vrchol zavřu
\item algoritmus začíná přidáním prvního vrcholu do fronty (to je ten vrchol, ze kterého graf prohledáváme)
\item algoritmus končí, jakmile je fronta prázdná
\item složitost $\Theta(n+m)$
\item BFS najde cestu s nejmenším počtem hran (z výchozího vrcholu do libovolného vrcholu grafu) - tedy lze použít k hledání nejkratší cesty v grafech s jednotkovými hranami 
\end{itemize}
\item prohledávání do hloubky (DFS)
\begin{itemize}
\item lze implementovat rekurzí nebo jako BFS se zásobníkem místo fronty
\item na začátku jsou všechny vrcholy neviděné
\item DFS(v):
\begin{itemize}
\item stav(v) $\leftarrow$ otevřený
\item pro vw $\in E$
\begin{itemize}
\item pokud stav(w) = neviděný
\begin{itemize}
\item DFS(w)
\end{itemize}
\end{itemize}
\item stav(v) $\leftarrow$ zavřený
\end{itemize}
\item DFS spouštíme z nějaké vrcholu - navštívíme všechny vrcholy z něj dosažitelné
\item opakované DFS - algoritmus nespouštíme pouze pro jeden určitý vrchol, ale pro všechny (neviděné) vrcholy grafu $\to$ projdeme všechny komponenty souvislosti
\begin{itemize}
\item stejného efektu lze docílit přidáním jednoho superzdroje, který bude připojen ke všem vrcholům grafu
\end{itemize}
\item složitost $\Theta(n+m)$
\end{itemize}
\end{itemize}
\item Topologické třídění orientovaných grafů
\begin{itemize}
\item topologické uspořádání grafu $G=(V,E)$ je lineární uspořádání $\leq$ na $V$ takové, že $\forall uv\in E: u\leq v$
\begin{itemize}
\item může jich být víc
\end{itemize}
\item konstrukce topologického uspořádání
\begin{itemize}
\item postupným odtrháváním zdrojů (vezmeme libovolný vrchol, jdeme proti šipkám do zdroje, tento zdroj umístíme do uspořádání a odtrhneme ho z grafu) - to je ale složité na efektivní implementaci
\item opakované DFS opouští (zavírá) vrcholy v pořadí opačném topologickému uspořádání
\end{itemize}
\end{itemize}
\item Nejkratší cesty v ohodnocených grafech (Dijkstrův a Bellmanův-Fordův algoritmus)
\begin{itemize}
\item Dijkstra
\begin{itemize}
\item hledání nejkratších cest z daného vrcholu do všech vrcholů grafu
\item je to v podstatě BFS s budíky (s prioritní frontou)
\item funguje pro kladné reálné délky hran
\item začínáme od nějakého vrcholu $v_0$ (otevřeme ho a jeho ohodnocení $h(v_0)$ nastavíme na nulu)
\begin{itemize}
\item ostatní vrcholy $v$ mají ohodnocení $h(v)=+\infty$
\end{itemize}
\item dokud existují nějaké otevřené vrcholy, vybereme z nich ten, jehož $h(v)$ je nejmenší (ten zavřeme) a všechny jeho následníky otevřeme a jejich $h(w)$ snížíme na $h(v)+l(v,w)$
\item ukládáme předchůdce vrcholů, aby se cesta dala rekonstruovat
\item pozorování: každý vrchol zavřeme pouze jednou
\item pokud nejmenší $h(v)$ hledáme pokaždé znova, tak to má složitost $\Theta(n^2)$
\item cena operací
\begin{itemize}
\item ExtractMin \dots{} $T_X$
\item Insert \dots{} $T_I$
\item Decrease \dots{} $T_D$
\end{itemize}
\item složitost Dijkstra je $O(n\cdot T_I+m\cdot T_D + n\cdot T_X)$
\begin{itemize}
\item vkládáme každý vrchol
\item decrease se provádí nejvýše jednou za každou hranu
\item extractujeme každý vrchol
\end{itemize}
\item při použití binární haldy jsou všechny tři operace $O(\log n)$, z čehož vyplývá složitost Dijkstra $O((n+m)\log n)$
\begin{itemize}
\item lepší (asymptoticky) je zvolit $d$-regulární haldu, kde $d=m/n$
\item ještě lepší je Fibonacciho halda
\end{itemize}
\end{itemize}
\item Bellman-Ford
\begin{itemize}
\item relaxační algoritmus (je to třída algoritmů, liší se tím, jak vybírají otevřené vrcholy - patří tam i Dijkstra, který vybírá ty nejlevnější)
\item vrcholy ukládá do fronty (takže jako Dijkstra, akorát nebere nejlevnější, ale nejstarší)
\item funguje pro grafy bez záporných cyklů
\item obecný relaxační algoritmus
\begin{itemize}
\item vstup: graf $G$, počáteční vrchol $v_0$
\item na začátku jsou všechny vrcholy nenalezené, jejich ohodnocení je nekonečné, jejich předchůdci jsou nedefinováni
\item stav počátečního vrcholu nastavíme na otevřený, jeho ohodnocení na nulu
\item dokud existují otevřené vrcholy
\begin{itemize}
\item vezmu nějaký otevřený vrchol $v$ (v Bellmanově-Fordově algoritmu beru ten nejstarší ve frontě)
\item pro všechny jeho následníky $w$ provedu relaxaci, tzn.:
\begin{itemize}
\item pokud $h(w)\gt h(v)+l(v,w)$
\begin{itemize}
\item $h(w)\leftarrow h(v)+l(v,w)$
\item stav(w) $\leftarrow$ otevřený
\item $P(w)\leftarrow v$
\end{itemize}
\item stav(v) $\leftarrow$ uzavřený
\end{itemize}
\end{itemize}
\item výstup: ohodnocení vrcholů $h$ a pole předchůdců
\end{itemize}
\item Bellmanův-Fordův algoritmus má fáze
\begin{itemize}
\item $F_0:=$ otevření $v_0$
\item $F_i:=$ zavírání vrcholů otevřených v $F_{i-1}$ a otevírání jejich následníků
\end{itemize}
\item invariant: na konci fáze $F_i$ odpovídá $h(v)$ délce nejkratšího $v_0v$-sledu o nejvýše $i$ hranách
\item z toho vyplývá složitost $\Theta(n\cdot m)$
\end{itemize}
\end{itemize}
\item Minimální kostra grafu (Jarníkův a Borůvkův algoritmus)
\begin{itemize}
\item lemma
\begin{itemize}
\item nechť $G$ je graf opatřený unikátními vahami, $R$ nějaký jeho elementární řez a $e$ nejlehčí hrana tohoto řezu
\item pak $e$ leží v každé minimální kostře grafu $G$
\end{itemize}
\item klasický Jarník
\begin{itemize}
\item hladový algoritmus
\item na začátku máme strom $T$, který obsahuje vrchol $v_0$ (libovolný) a žádné hrany
\item vezmeme nejlehčí hranu vedoucí mezi stromem $T$ a zbytkem grafu - tu přidáme do stromu
\begin{itemize}
\item opakujeme, dokud takové hrany existují
\end{itemize}
\item složitost $O(n\cdot m)$
\end{itemize}
\item Jarník podle Dijkstry
\begin{itemize}
\item udržuju si haldu aktivních hran - to jsou hrany v aktuálním řezu
\item do stromu $T$ vždycky přidávám minimum z haldy
\item když zjistím, že mám dvě aktivní hrany, které vedou do jednoho vrcholu mimo $T$, tak tu těžší zahodím
\item s haldou složitost $O(m\log n)$
\begin{itemize}
\item protože $n\log n$ se díky souvislosti grafu vejde do $m\log n$
\end{itemize}
\end{itemize}
\item Borůvkův algoritmus
\begin{itemize}
\item paralelní verze Jarníkova algoritmu
\item na začátku $T$ obsahuje izolované vrcholy grafu
\item dokud $T$ není souvislý:
\begin{itemize}
\item rozložíme $T$ na komponenty souvislosti
\item pro každou komponentu najdeme nejlehčí hranu mezi komponentou a zbytkem vrcholů a přidáme ji do $T$
\end{itemize}
\item složitost
\begin{itemize}
\item rozklad na komponenty v $O(n+m)$, ale taky $O(m)$ díky souvislosti
\item nalezení nejlehčích hran v $O(m)$
\item algoritmus se zastaví po nejvýše $\lfloor\log n\rfloor$ iteracích
\item z toho plyne složitost $O(m\log n)$
\end{itemize}
\end{itemize}
\end{itemize}
\item Toky v sítích (Ford-Fulkerson algoritmus)
\begin{itemize}
\item cesta je nenasycená $\equiv$ žádná její hrana není nasycená / všechny mají kladné rezervy
\item algoritmus
\begin{itemize}
\item iterujeme, dokud existuje nenasycená cesta ze zdroje do stoku
\item spočítáme rezervu celé cesty (minimum přes rezervy hran cesty)
\item pro každou hranu upravíme tok - v~protisměru odečteme co nejvíc, zbytek přičteme po směru
\end{itemize}
\item pro celočíselné kapacity vrátí celočíselný tok
\item racionální kapacity převedeme na celočíselné
\item pro iracionální kapacity se může rozbít
\item lemma: pokud se algoritmus zastaví, vydá maximální tok
\item věta: pro každou síť s racionálními kapacitami se Fordův-Fulkersonův algoritmus zastaví a vydá maximální tok a minimální řez
\begin{itemize}
\item vzhledem k operacím, které algoritmus provádí, nemůže z celých čísel vytvořit necelá
\end{itemize}
\item důsledek: velikost maximálního toku je rovna kapacitě minimálního řezu
\item důsledek: síť s celočíselnými kapacitami má aspoň jeden z maximálních toků celočíselný a Fordův-Fulkersonův algoritmus takový tok najde
\end{itemize}
\end{itemize}
\subsection{3. Programovací jazyky}
\begin{itemize}
\item Abstrakce, zapouzdření, polymorfismus - související konstrukty programovacích jazyků (třídy, rozhraní, metody, datové položky, dědičnost, viditelnost)
\begin{itemize}
\item v C\# i v C++ jsou třídy (class) i struktury (struct)
\begin{itemize}
\item deklarace musí být v C++ zakončena středníkem
\item v C++ to neovlivňuje způsob uložení a předávání, v C\# jsou struktury hodnotového typu
\end{itemize}
\item rozhraní (interface) v C\#
\begin{itemize}
\item konvence psát na začátek názvu \texttt{I}
\item dovnitř se píšou metody / method-like věci (třeba props), nesmí tam být fields
\item jeden řádek může být např. \texttt{public \allowbreak{}int \allowbreak{}m(\allowbreak{}int \allowbreak{}a,\allowbreak{} \allowbreak{}int \allowbreak{}b)\allowbreak{};\allowbreak{}}
\item dříve muselo být všechno public, dneska je dobré tam to public explicitně psát (i když defaultně je pořád všechno public)
\item názvy parametrů metod v interfacu slouží k dokumentaci
\item nemůže existovat instance interfacu, pouze proměnná s typem interfacu (ta může být nullable)
\item \texttt{class \allowbreak{}Kachna \allowbreak{}: \allowbreak{}IPostovniZvire \allowbreak{}\{...\}\allowbreak{}}
\item \texttt{IPostovniZvire \allowbreak{}zvire1 \allowbreak{}= \allowbreak{}new \allowbreak{}Kachna(\allowbreak{})\allowbreak{};\allowbreak{}}
\item třída může implementovat víc interfaců
\item rozhraní může dědit od jiného rozhraní
\end{itemize}
\item v C++ můžeme rozhraní definovat pomocí dědičnosti (obvykle virtuální)
\item v C++ můžu při dědění definovat viditelnost
\item viditelnost v C\# - některá klíčová slova
\begin{itemize}
\item public - přístup není omezen
\item private - přístup je omezen na kód v daném typu
\begin{itemize}
\item C\# podporuje vnořené typy - takže kód ve vnořeném typu má taky přístup k private položkám typu, v němž je vnořen
\end{itemize}
\item protected - přístup je omezen na daný typ a typy, které z něj dědí
\end{itemize}
\item výchozí viditelnost - private ve třídě, public v interfacu
\item C\#: abstract method vs. interface
\begin{itemize}
\item vynucení kontraktu - interface ho vynucuje hned
\item veřejný kontrakt - abstraktní metody můžou být protected
\item slib rozšiřitelnosti - u virtuální metody je to opt-out (pomocí sealed), u interfaců opt-in (musím znova říct, že implementuju interface)
\item data - u interfaců nelze (kvůli vícenásobné dědičnosti)
\item vícenásobná dědičnost - nelze u abstraktních metod
\end{itemize}
\end{itemize}
\item (Dynamický) polymorfismus, statické a dynamické typování
\begin{itemize}
\item statické typování
\begin{itemize}
\item proměnné mají pevně dané typy známé při kompilaci
\item je jasné, jaké metody podporují
\item dobře se provádí statická analýza kódu, lze hlásit chybná volání metod apod.
\end{itemize}
\item dynamické typování
\begin{itemize}
\item typy proměnných nejsou pevně definované, určují se za běhu
\item volání chybné metody selže až za běhu
\item tzv. duck typing
\item v C\# lze aktivovat klíčovým slovem \texttt{dynamic}
\item za jistou formu dynamického typování lze považovat i používání obecných rozhraní nebo rodičovských typů (např. \texttt{object})
\end{itemize}
\item statický polymorfismus se implementuje tak, že přetížíme metodu nebo operátor nebo že zakryjeme rodičovskou metodu
\item dynamický polymorfismus funguje pomocí virtuálních metod
\item při volání interfacové metody je důležité, která třída ji implementuje
\begin{itemize}
\item koncept virtuálních metod s interfacy přímo nesouvisí, ale požadavek interfacu můžeme uspokojit virtuální metodou
\item poznámka: interfacová metoda se dá implementovat explicitně
\end{itemize}
\end{itemize}
\item Vícenásobná dědičnost a její problémy (vícenásobná a virtuální dědičnost v C++, interfaces v C\#)
\begin{itemize}
\item diamantový problém - pokud mají moji rodiče společného rodiče a já volám jeho metodu, kterou oba implementují, jaká metoda se má volat?
\end{itemize}
\item Číselné a výčtové typy
\begin{itemize}
\item enum se v C\# definuje třeba takto: \texttt{enum \allowbreak{}Season \allowbreak{}\{ \allowbreak{}Spring,\allowbreak{} \allowbreak{}Summer,\allowbreak{} \allowbreak{}Autumn,\allowbreak{} \allowbreak{}Winter \allowbreak{}\}\allowbreak{}}
\begin{itemize}
\item když za Season napíšu \texttt{: \allowbreak{}ushort}, tak se místo intu použije ushort
\item můžu nadefinovat hodnoty jednotlivých možností
\end{itemize}
\item v C\# jsou tyto číselné typy: byte, sbyte, decimal, double, float, int, uint, nint, nuint, long, ulong, short, ushort
\begin{itemize}
\item nint a nuint mají nativní velikost (64 nebo 32 bitů, podle platformy)
\end{itemize}
\end{itemize}
\item Ukazatele a reference v C++
\begin{itemize}
\item smart pointery
\item raw pointery
\item reference (const, rvalue)
\end{itemize}
\item Hodnotové a referenční typy v C\#
\begin{itemize}
\item základní dělení všech typů
\begin{itemize}
\item pointery - jsou ošklivé, nebudeme se o nich bavit
\item hodnotové typy - alokované na místě (s výjimkami)
\begin{itemize}
\item enums
\item structures
\begin{itemize}
\item simple types (Int32, Int64, Double, Boolean, Char, \dots{})
\item nullables
\item user defined structures (struct)
\end{itemize}
\end{itemize}
\item referenční typy - alokované na spravované (managed) haldě
\begin{itemize}
\item classes (e.g. strings)
\item interfaces
\item arrays
\item delegates
\end{itemize}
\end{itemize}
\item halda je garbage-collectovaná (GC), funguje chytřeji než v Pythonu, není potřeba reference counter, ale používá se graf dosažitelnosti
\item overhead u každého objektu na haldě (má typicky 16 B)
\begin{itemize}
\item syncblock - kvůli práci s více vlákny (u jednovláknových programů zbytečný), má 8 B, skládá se ze dvou částí:
\begin{itemize}
\item lock
\begin{itemize}
\item owner thread + counter (kvůli rekurzivnímu zamykání)
\item waiting threads list - nějaký férový seznam (nebo fronta, na tom nesejde)
\end{itemize}
\item condition variable
\begin{itemize}
\item sleeping threads list
\end{itemize}
\end{itemize}
\item pointer na typ (zjednodušeně řečeno)
\begin{itemize}
\item třída System.Type, má instance na GC haldě
\item každý datový typ odpovídá jedné instanci třídy
\end{itemize}
\item na referenčních proměnných jde volat \texttt{.\allowbreak{}GetType(\allowbreak{})\allowbreak{}}
\end{itemize}
\item co může být v proměnné
\begin{itemize}
\item hodnota (u hodnotových typů)
\begin{itemize}
\item velikost odpovídá velikosti datového typu
\item u struktur lze použít \texttt{new}, které nic nealokuje, ale zavolá konstruktor
\end{itemize}
\item reference (u referenčních typů)
\begin{itemize}
\item je tam uložená adresa, takže velikost odpovídá velikosti adres
\item adresa vede na GC haldu
\item adresa ukazuje na celý objekt
\item explicitní alokace pomocí \texttt{new}
\item tu adresu nelze zjistit (zčásti proto, že GC objekty na haldě někdy přesouvá, proto se adresa může měnit)
\end{itemize}
\end{itemize}
\end{itemize}
\item Pokročilé prvky syntaxe C\#
\begin{itemize}
\item vícerozměrná pole
\begin{itemize}
\item zubatá (jagged)
\begin{itemize}
\item pole polí
\begin{itemize}
\item \texttt{int[][] \allowbreak{}a \allowbreak{}= \allowbreak{}new \allowbreak{}int[2][];\allowbreak{}}
\item \texttt{a[0] \allowbreak{}= \allowbreak{}new \allowbreak{}int[3];\allowbreak{}}
\item \texttt{a[1] \allowbreak{}= \allowbreak{}new \allowbreak{}int[4];\allowbreak{}}
\item \texttt{int \allowbreak{}x \allowbreak{}= \allowbreak{}a[0][1];\allowbreak{}}
\end{itemize}
\item jako v Javě
\item nevýhoda: každé pole má svůj vlastní overhead
\end{itemize}
\item obdélníková (rectangular)
\begin{itemize}
\item mezi hranaté závorky napíšu jednu nebo více čárek
\item \texttt{int[,\allowbreak{}] \allowbreak{}a \allowbreak{}= \allowbreak{}new \allowbreak{}int[2,\allowbreak{} \allowbreak{}3];\allowbreak{}}
\item \texttt{int \allowbreak{}x \allowbreak{}= \allowbreak{}a[0,\allowbreak{} \allowbreak{}1];\allowbreak{}}
\item jako v C/C++
\end{itemize}
\item jagged jsou obecně rychlejší (cca 2$\times$)
\begin{itemize}
\item ale pokud už s každou naindexovanou položkou pole budu provádět nějaký výpočet, tak tam nebude velký rozdíl v rychlosti
\end{itemize}
\item obdélníková jsou paměťově úspornější (v určitých situacích), lépe se zapisují
\end{itemize}
\item indexer je vlastně speciální property
\begin{itemize}
\item \texttt{public \allowbreak{}T \allowbreak{}this[int \allowbreak{}i] \allowbreak{}\{ \allowbreak{}get \allowbreak{}\{ \allowbreak{}return \allowbreak{}arr[i];\allowbreak{} \allowbreak{}\}\allowbreak{} \allowbreak{}set \allowbreak{}\{ \allowbreak{}arr[i] \allowbreak{}= \allowbreak{}value;\allowbreak{} \allowbreak{}\}\allowbreak{} \allowbreak{}\}\allowbreak{} \allowbreak{}\}\allowbreak{}}
\end{itemize}
\item nullable value types mají vlastnost \texttt{bool \allowbreak{}HasValue} a taky \texttt{T \allowbreak{}Value}
\item přetěžování operátorů
\begin{itemize}
\item \texttt{public \allowbreak{}static \allowbreak{}Fraction \allowbreak{}operator \allowbreak{}*(\allowbreak{}Fraction \allowbreak{}left,\allowbreak{} \allowbreak{}Fraction \allowbreak{}right)\allowbreak{} \allowbreak{}=\textgreater{}\allowbreak{} \allowbreak{}new \allowbreak{}Fraction(\allowbreak{}left.\allowbreak{}numerator \allowbreak{}* \allowbreak{}right.\allowbreak{}numerator,\allowbreak{} \allowbreak{}left.\allowbreak{}denominator \allowbreak{}* \allowbreak{}right.\allowbreak{}denominator)\allowbreak{};\allowbreak{}}
\end{itemize}
\item přetížení konverze
\begin{itemize}
\item \texttt{public \allowbreak{}static \allowbreak{}implicit \allowbreak{}operator \allowbreak{}byte(\allowbreak{}Digit \allowbreak{}d)\allowbreak{} \allowbreak{}=\textgreater{}\allowbreak{} \allowbreak{}d.\allowbreak{}digit;\allowbreak{}}
\item \texttt{public \allowbreak{}static \allowbreak{}explicit \allowbreak{}operator \allowbreak{}Digit(\allowbreak{}byte \allowbreak{}b)\allowbreak{} \allowbreak{}=\textgreater{}\allowbreak{} \allowbreak{}new \allowbreak{}Digit(\allowbreak{}b)\allowbreak{};\allowbreak{}}
\end{itemize}
\end{itemize}
\item Reference, imutabilní typy a boxing v C\#
\begin{itemize}
\item boxing
\begin{itemize}
\item každý hodnotový typ je potomkem objectu
\item tedy do proměnné typu object musí jít přiřadit proměnnou hodnotového typu
\item při takovém přiřazení se provádí boxing - je to implicitní alokace na GC haldě, alokuje se tam instance hodnotového typu (bude tam standardní overhead)
\begin{itemize}
\item pozor: v Javě jsou dva různé typy, int (hodnotový) a Integer (referenční) - v C\# je to ten samý typ System.Int32
\end{itemize}
\item co můžu dělat s proměnnou typu object?
\begin{itemize}
\item můžu volat ToString()
\end{itemize}
\item unboxing - hodnotový typ alokovaný na haldě uložený referencí se uloží jako hodnota
\item není vhodné používat object pro slabé typování - vlastně takhle implementujeme duck typing
\item když potřebujeme hodnotový typ předávat referencí, tak object nedává smysl, vhodnější je použít jednoduchou třídu
\end{itemize}
\item imutabilní typy
\begin{itemize}
\item do fieldu označeného jako readonly lze zapisovat pouze v konstruktoru (a při inicializaci apod.)
\item případně u vlastností můžeme schovat setter - pak budou taky readonly
\item strukturu můžeme celou označit jako readonly
\end{itemize}
\item reference
\begin{itemize}
\item hodí se pro snazší manipulaci s hodnotovými proměnnými (abychom je všude nemuseli předávat hodnotou a abychom mohli z funkce vracet víc hodnot než jednu)
\begin{itemize}
\item u referenčních typů obvykle nedává smysl používat tracking referenci
\item výjimečným případem, kdy to smysl dává, je třeba zvětšení pole - vyrobím nové pole, položky překopíruju a do tracking reference přiřadím nové pole
\end{itemize}
\item používají se tzv. tracking reference
\item klíčové slovo \texttt{ref} - používá se u parametrů funkcí, uvádí se dvakrát, v~hlavičce funkce i při volání
\item když používám \texttt{ref}, tak proměnná musí být inicializovaná (aby se z ní dalo číst)
\item proto se u výstupních parametrů funkcí používá klíčové slovo \texttt{out} - na úrovni CIL kódu je to to samé, ale nevyžaduje to, aby proměnné byly inicializované
\begin{itemize}
\item uvnitř funkce s výstupními parametry se s nimi zachází jako s~neinicializovanými (dá se z nich číst až po prvním zápisu)
\item před standardním ukončením funkce s výstupními parametry se do každého z~nich musí alespoň jednou zapsat
\item pokud funkce skončí výjimkou, tak se do výstupních parametrů zapsat nemusí, ale to není problém, protože catch blok nemůže spoléhat na inicializaci v try bloku
\item pokud proměnnou deklaruju nad try/catch, v try bloku ji inicializuji a v catch bloku vyhodím výjimku dál (pomocí \texttt{throw;\allowbreak{}}), tak pod try/catch můžu proměnnou považovat za inicializovanou
\item pokud napíšu \texttt{if \allowbreak{}(\allowbreak{}int.\allowbreak{}Parse(\allowbreak{}s,\allowbreak{} \allowbreak{}out \allowbreak{}int \allowbreak{}x)\allowbreak{})\allowbreak{}}, tak se proměnná deklaruje před ifem
\end{itemize}
\item klíčové slovo \texttt{in} \dots{} proměnná uvnitř funkce funguje jako readonly (podobně celá struktura funguje jako readonly)
\begin{itemize}
\item při volání funkce se dá volat i hodnotově bez klíčového slova
\item \texttt{in} má smysl u výkonnostních optimalizací hodnotových typů (typicky u počítačových her)
\end{itemize}
\item místo \texttt{in} se dá použít taky \texttt{ref \allowbreak{}readonly} - je to něco velmi podobného
\item tracking reference se dají používat i jinde (nejen jako parametry funkcí) - je lepší je moc nepoužívat
\end{itemize}
\end{itemize}
\item Šablony (templates) a statický polymorfismus v C++
\begin{itemize}
\item např. \texttt{template\textless{}typename \allowbreak{}T\textgreater{}\allowbreak{} \allowbreak{}const \allowbreak{}T\& \allowbreak{}max(\allowbreak{}const \allowbreak{}T\& \allowbreak{}a,\allowbreak{} \allowbreak{}const \allowbreak{}T\& \allowbreak{}b)\allowbreak{} \allowbreak{}\{ \allowbreak{}return \allowbreak{}a \allowbreak{}\textless{} \allowbreak{}b \allowbreak{}? \allowbreak{}b \allowbreak{}: \allowbreak{}a;\allowbreak{} \allowbreak{}\}\allowbreak{}}
\begin{itemize}
\item templatovaná funkce
\item v C\# třeba pomocí generické metody a omezení na IComparable
\end{itemize}
\end{itemize}
\item Generické typy v C\# (bez omezení typových parametrů)
\begin{itemize}
\item \texttt{class \allowbreak{}GenericList\textless{}T\textgreater{}\allowbreak{} \allowbreak{}\{ \allowbreak{}\}\allowbreak{}}
\item při vytvoření instance je potřeba uvést typový parametr (při volání generických metod to není potřeba)
\end{itemize}
\item Typy reprezentující funkce v C++, C\#
\begin{itemize}
\item v C\# se používají delegáti
\begin{itemize}
\item pomocí klíčového slova \texttt{delegate} definujeme nový delegátní typ: \texttt{public \allowbreak{}delegate \allowbreak{}void \allowbreak{}Callback(\allowbreak{}string \allowbreak{}message)\allowbreak{};\allowbreak{}}
\item za předpokladu, že někde existuje \texttt{DelegateMethod} s jedním parametrem typu string, která vrací void, tak můžeme provést přiřazení \texttt{Callback \allowbreak{}handler \allowbreak{}= \allowbreak{}DelegateMethod;\allowbreak{}}
\begin{itemize}
\item kdybychom chtěli vynutit vytvoření nové instance delegáta, tak bychom použili syntaxi \texttt{Callback \allowbreak{}handler \allowbreak{}= \allowbreak{}new \allowbreak{}Callback(\allowbreak{}DelegateMethod)\allowbreak{};\allowbreak{}}
\end{itemize}
\item pak ji můžeme použít: \texttt{handler(\allowbreak{}"Hello \allowbreak{}World")\allowbreak{};\allowbreak{}}
\item k viditelnosti: můžu mít public metodu, která vrací delegáta, který ukazuje na private statickou metodu
\item delegáti můžou ukazovat i na instanční metody
\begin{itemize}
\item v delegátu jsou dva ukazatele - na funkci a na \texttt{this}
\begin{itemize}
\item u statických metod je tam null
\end{itemize}
\item stále platí, že delegáti jsou immutable - to konkrétní \texttt{this} je tam navždy
\item u struktur - do \texttt{this} se zkopíruje struktura (zaboxuje se)
\end{itemize}
\item delegáti můžou být generičtí
\item u delegátů se dává suffix Delegate
\item občas mi jde čistě o parametry - je mi úplně jedno, co ten delegát dělá
\begin{itemize}
\item bylo by zbytečné pro každou takovou metodu deklarovat delegáta
\item proto existuje spousta předdefinovaných delegátů
\begin{itemize}
\item \texttt{Action\textless{}...\textgreater{}\allowbreak{}(\allowbreak{}...)\allowbreak{} \allowbreak{}$\to$ \allowbreak{}void}
\item \texttt{Func\textless{}...,\allowbreak{} \allowbreak{}TResult\textgreater{}\allowbreak{}(\allowbreak{}...)\allowbreak{} \allowbreak{}$\to$ \allowbreak{}TResult}
\item \texttt{Predicate\textless{}T\textgreater{}\allowbreak{}(\allowbreak{}T \allowbreak{}obj)\allowbreak{} \allowbreak{}$\to$ \allowbreak{}bool}
\end{itemize}
\item pozor na přehlednost - typicky je lepší používat silně typované delegáty
\end{itemize}
\item \texttt{var} funguje i pro delegáty
\end{itemize}
\item v C++ je víc způsobů, jak pracovat s funkcemi
\begin{itemize}
\item templates
\item function pointers
\item \texttt{std::function}
\end{itemize}
\end{itemize}
\item Lambda funkce a funkcionální rozhraní
\begin{itemize}
\item lambda funkce v C\# se píšou jako \texttt{(\allowbreak{}input-parameters)\allowbreak{} \allowbreak{}=\textgreater{}\allowbreak{} \allowbreak{}expression} nebo \texttt{(\allowbreak{}input-parameters)\allowbreak{} \allowbreak{}=\textgreater{}\allowbreak{} \allowbreak{}\{ \allowbreak{}\textless{}sequence-of-statements\textgreater{}\allowbreak{} \allowbreak{}\}\allowbreak{}}
\item můžou být statické
\begin{itemize}
\item např. \texttt{static \allowbreak{}x \allowbreak{}=\textgreater{}\allowbreak{} \allowbreak{}x \allowbreak{}+ \allowbreak{}1}
\end{itemize}
\item pokud nejsou statické, můžou zachytávat proměnné z kontextu (provádí se zachycení podobné capture by reference v C++)
\item lambda funkce jsou přiřaditelné do proměnných typu delegát
\end{itemize}
\item Správa životního cyklu zdrojů v případě výskytu chyb - RAII v C++, using v C\#, obsluha a propagace výjimek
\begin{itemize}
\item RAII \dots{} Resource Acquisition is Initialization
\begin{itemize}
\item zdroje, které třída používá, by se měly pořídit v konstruktoru a uvolnit v destruktoru, aby to bylo bezpečné
\item lock \& unlock $\to$ lock\_guard
\item new \& delete $\to$ smart pointers
\item pokud chceme ohraničit část kódu, kde má být zámek zamčený (nebo soubor otevřený), tak ji uzavřeme do složených závorek
\end{itemize}
\item v C\# pomocí usingu
\begin{itemize}
\item \texttt{using \allowbreak{}(\allowbreak{}var \allowbreak{}f1 \allowbreak{}= \allowbreak{}new \allowbreak{}FileStream(\allowbreak{}".\allowbreak{}.\allowbreak{}.\allowbreak{}")\allowbreak{})\allowbreak{} \allowbreak{}\{ \allowbreak{}.\allowbreak{}.\allowbreak{}.\allowbreak{} \allowbreak{}\}\allowbreak{}}
\item je to syntaktická zkratka pro try \& finally, přičemž ve finally bloku se volá Dispose na hodnotě výrazu v závorce
\end{itemize}
\item výjimky
\begin{itemize}
\item vyhazují se pomocí \texttt{throw \allowbreak{}new \allowbreak{}Ex(\allowbreak{})\allowbreak{};\allowbreak{}}, kde \texttt{Ex} je nějaký typ výjimky (konstruktor může mít parametr - třeba nějakou zprávu)
\item zachycují/zpracovávají se pomocí try-catch-finally bloku
\item catch bloků může být víc pro různé typy výjimek
\item v catch bloku můžeme tu stejnou výjimku vyhodit znova pomocí \texttt{throw;\allowbreak{}}
\end{itemize}
\end{itemize}
\item Alokace (alokace statická, na zásobníku, na haldě)
\begin{itemize}
\item statická alokace se provádí už během kompilace, týká se proměnných, které mají být k dispozici během celého běhu programu
\item alokace na zásobník se děje při volání funkcí - lokální proměnné se takto alokují a po doběhnutí funkce se zase smažou
\begin{itemize}
\item je to rychlejší než alokace na haldě, ale na (volacím) zásobníku je méně místa než na haldě
\item opravdu to funguje jako zásobník (FIFO)
\item velikost alokované paměti je známá během kompilace
\end{itemize}
\item dynamickou alokaci na haldě použijeme v situacích, kdy nám první dva přístupy nestačí
\begin{itemize}
\item potřebujeme sdílet objekty nějak napříč
\item potřebujeme polymorfismus a nestačí nám reference (v C++)
\item potřebujeme alokovat hodně paměti
\end{itemize}
\item v C++ je nutné dynamicky alokovanou paměť ručně dealokovat (smart pointers tohle řeší)
\item v C\# jsou všechny instance referenčních typů (a jejich fieldy) na haldě
\end{itemize}
\item Inicializace (konstruktory, volání zděděných konstruktorů)
\begin{itemize}
\item v C\# je konstruktor označen viditelností a názvem třídy, volá se s new
\item volání rodičovského konstruktoru zajišťuje klíčové slovo base
\begin{itemize}
\item \texttt{public \allowbreak{}B(\allowbreak{}params1)\allowbreak{} \allowbreak{}: \allowbreak{}base(\allowbreak{}params2)\allowbreak{} \allowbreak{}\{}
\item není to optional, rodičovský konstruktor se volá pokaždé (defaultně bez parametrů)
\end{itemize}
\item konstruktor předka se vždy volá před konstruktorem potomka
\item v C++ je syntaxe velmi podobná, akorát místo \texttt{base} píšeme přímo název rodičovské třídy
\end{itemize}
\item Destrukce (destruktory, finalizátory)
\begin{itemize}
\item v C\# se finalizátory používají k explicitnímu uvolňování unmanaged zdrojů (ve specifických use-casech)
\item v C++ se destruktory používají častěji, nejtypičtější to je, pokud potřebujeme mít pod kontrolou vytváření kopií objektu apod. (viz rule of five)
\item rodičovské třídy v C++ by měly mít virtuální destruktor (stačí prázdný)
\end{itemize}
\item Explicitní uvolňování objektů, reference counting, garbage collector
\begin{itemize}
\item v C++ se na haldě alokuje pomocí new, pak je nutné objekt smazat pomocí delete (právě jednou)
\begin{itemize}
\item alternativou jsou smart pointers
\begin{itemize}
\item unique\_ptr dealokuje, jakmile dochází se smazání pointeru (proměnná vypadne ze scopu)
\item shared\_ptr umožňuje ukazovat na jeden objekt z více míst, počítá reference - dealokuje, jakmile je referencí nula
\begin{itemize}
\item může být problém s cyklickými referencemi, proto se někdy používá weak\_ptr
\end{itemize}
\end{itemize}
\item v C se k podobnému účelu používá malloc a free
\end{itemize}
\item v C\# je halda garbage collectovaná
\begin{itemize}
\item garbage collector se pouští podle toho, zda je potřeba uvolnit paměť
\item fungují tam tři generace (vrstvy) objektů
\item vytváří se graf objektů, aby bylo jasné, které jsou používané a které je možné smazat
\item po smazání se zbylé objekty na haldě přeskupují (tzv. heap compacting)
\item GC má dva módy - workstation a server (v server módu běží GC na více vláknech)
\end{itemize}
\end{itemize}
\item Reprezentace vláken v programovacích jazycích
\begin{itemize}
\item v C\# je ve jmenném prostoru System.Threading třída Thread, která obvykle reprezentuje jedno vlákno
\item ale vyrobení vlákna je drahé, takže může být užitečnější použít ThreadPool
\end{itemize}
\item Specifikace funkce vykonávané vláknem a základní operace nad vlákny
\begin{itemize}
\item v C\# můžeme vytvořit instanci třídy Thread
\item instanci předáváme delegáta, ten může být bez parametrů nebo s jedním parametrem typu object
\item vlákno spustíme pomocí volání Start()
\item pokud v nějaký moment chceme počkat, než doběhne, můžeme použít Join()
\end{itemize}
\item Časově závislé chyby (race condition) a mechanismy pro synchronizaci vláken
\begin{itemize}
\item race condition nastává, když je výsledek operace závislý na pořadí, v jakém se provedou události, jejichž pořadí nemáme pod kontrolou
\item nejjednodušší řešení je mutual exclusion - k tomu se používá zámek
\begin{itemize}
\item \texttt{Monitor.\allowbreak{}Enter(\allowbreak{}obj)\allowbreak{}} zamkne zámek na daném objektu (nebo pasivně čeká, dokud se neodemkne)
\item \texttt{Monitor.\allowbreak{}Exit(\allowbreak{}obj)\allowbreak{}} odemkne zámek
\item existuje syntaktická zkratka \texttt{lock(\allowbreak{}obj)\allowbreak{} \allowbreak{}\{ \allowbreak{}... \allowbreak{}\}\allowbreak{}}, na začátku bloku se volá Enter, na konci se volá Exit
\item zámek si v syncblocku ukládá referenci na vlákno, které ho vlastní (jinak je tam null)
\item pokud ho zamkneme ze stejného vlákna několikrát, musíme ho několikrát odemknout (k tomu má počítadlo)
\end{itemize}
\item syncblock má každá referenční proměnná
\begin{itemize}
\item co budeme zamykat?
\item chceme mít jistotu, že nám to nezamkne nikdo jiný proti naší vůli
\item tedy pokud zamykáme například nějakou naši vlastní implementaci seznamu, dává smysl, aby ta třída měla soukromý field typu object, který bude sloužit jenom jako zámek
\item neměli bychom používat \texttt{lock \allowbreak{}(\allowbreak{}this)\allowbreak{}}
\end{itemize}
\item součástí syncblocku je taky seznam čekajících vláken určený pro tzv. condition variable
\begin{itemize}
\item vlákna čekají na to, až se s objektem něco stane
\begin{itemize}
\item třeba až se ve frontě objeví nějaké položky
\item tuhle podmínku je vhodné někam poznamenat, třeba do komentáře, je to ta \quotedblbase{}condition variable\textquotedblleft{}
\end{itemize}
\item když vlákno vidí, že je fronta prázdná, tak se pomocí \texttt{Monitor.\allowbreak{}Wait(\allowbreak{}obj)\allowbreak{}} zařadí do seznamu
\item jakmile se do fronty zařadí další položka, fronta může čekající vlákna notifikovat pomocí \texttt{Monitor.\allowbreak{}Pulse(\allowbreak{}obj)\allowbreak{}} nebo \texttt{Monitor.\allowbreak{}PulseAll(\allowbreak{}obj)\allowbreak{}}
\item aby se dal použít Wait, Pulse nebo PulseAll, musí být obalený v lock bloku (pro stejný objekt) - jinak to vyhodí výjimku
\end{itemize}
\item dalším synchronizačním primitivem je semafor, používá se k tomu, aby s objektem najednou mohlo pracovat jen omezené množství vláken
\end{itemize}
\item Implementace dynamického polymorfismu (tabulka virtuálních metod)
\begin{itemize}
\item každý typ s virtuálními metodami má tabulku virtuálních metod (VMT), ta se při dědění kopíruje (a prodlužuje), u overridnutých virtuálních metod se změní odkaz na implementaci
\item takže to, která implementace se reálně zavolá, závisí na třídě, jejíž instanci jsme vytvořili (nehledě na typ)
\end{itemize}
\item Nativní a interpretovaný běh, reprezentace programu, bytecode, interpret jazyka, just-in-time (JIT) a ahead-of-time (AOT) překlad
\begin{itemize}
\item nativní běh - zdrojový kód se přeloží a vznikne program, který lze přímo spustit
\item interpretovaný běh - zdrojový kód je zpracováván interpretem jazyka
\item C++
\begin{itemize}
\item zdrojáky .cpp, jeden po druhém je překládáme, vznikají object fily .obj (nebo .o), tam je přímo strojový kód
\begin{itemize}
\item ty se linkerem zpracují do jednoho .exe souboru
\end{itemize}
\item strojový kód typicky cílí na konkrétní platformu
\item kdybych chtěl, aby program běžel na jiné platformě, vezmu zdrojový kód a celý ho znova přeložím
\item 64bitové Windows podporují 32bitové programy
\item uživatelé musí vědět, jakou verzi programu si mají vybrat
\item knihovny
\begin{itemize}
\item překladač programu používá hlavičkové soubory knihoven
\end{itemize}
\end{itemize}
\item Apple
\begin{itemize}
\item Apple 6502 $\to$ Motorola 68000 (32bit) $\to$ IBM PowerPC (32bit) $\to$ Intel x86 (32bit) $\to$ Intel x64 (64bit) $\to$ Apple M1/M2 (ARM64, 64bit)
\item vymysleli fat binary (universal executable) - jeden spustitelný soubor, v něm je více verzí programu najednou
\item řeší to problém zpětné kompatibility, ne dopředné
\end{itemize}
\item C\#
\begin{itemize}
\item soubory .cs
\item .csproj, předává se build systému dotnetu, ten pustí C\# compiler csc.exe (dnes csc.dll)
\item vypadne z toho executable, uvnitř je CIL kód (common intermediate language)
\item ke spuštění programu potřebujeme CLR (common language runtime), obvykle je tam JIT (just-in-time) překladač, ten zajišťuje překlad a spuštění pro konkrétní platformu
\begin{itemize}
\item alternativou bylo vykonávat CIL instrukce pomocí běhové podpory, ale to je hrozně pomalé
\end{itemize}
\item CLR \dots{} virtual machine (VM)
\item CIL kód \dots{} managed/řízený kód
\item dotnet executable soubor \dots{} assembly
\item Java to má podobně
\begin{itemize}
\item Java intermediate language = Java bytecode
\end{itemize}
\item všechny programovací jazyky dotnetu se překládají do CIL kódu
\item optimalizace dělá až JIT
\begin{itemize}
\item ale má na to omezený čas, aby se program spustil dostatečně rychle
\end{itemize}
\item JIT používá překlad po metodách
\item ale dá se použít AOT (ahead of time) překlad - takže pak v .exe souboru je nejen assembly (ten se použije pro nepodporované platformy), ale taky přímo strojový kód, který tak může být efektivnější
\begin{itemize}
\item některé věci tam nefungují
\end{itemize}
\end{itemize}
\end{itemize}
\item Proces sestavení programu, oddělený překlad, linkování, staticky a dynamicky linkované knihovny
\begin{itemize}
\item sestavení programu
\begin{itemize}
\item preprocesor
\item kompilátor (překladač)
\item assembler
\item linker
\end{itemize}
\item knihovna - statická nebo dynamická sada binárek
\item linking - spojení binárek do jedné
\item loader - načte program do paměti
\item knihovna se linkuje do souboru .lib (statická knihovna)
\begin{itemize}
\item lidi už pak používají přímo hlavičkový soubor a statickou zkompilovanou knihovnu
\item pokud je knihovna statická, použije ji linker, pokud je dynamická, použije ji až loader
\end{itemize}
\item v C\# se o tohle všechno stará CLR (tedy .NET runtime), viz výše
\begin{itemize}
\item používají se dynamicky linkované knihovny (DLL)
\end{itemize}
\end{itemize}
\item Běhové prostředí procesu a vazba na operační systém
\begin{itemize}
\item program vs. proces (proces je spuštěný program, entita operačního systému)
\item operační systém zajišťuje, aby měl proces k dispozici virtuální paměť, aby mohl přistupovat k souborům na disku, k periferiím apod.
\item v C\# je běhovým prostředím CLR
\end{itemize}
\item Návrhové vzory
\begin{itemize}
\item decorator
\begin{itemize}
\item chceme nějaký objekt obalit, přidat mu funkce
\item dekorátor obvykle přijímá původní objekt v konstruktoru
\item má stejné rozhraní jako původní objekt
\end{itemize}
\item adapter
\begin{itemize}
\item převádí rozhraní jedné třídy na rozhraní jiné třídy
\item abychom instanci jiné třídy mohli použít tam, kde bychom normálně potřebovali instanci té původní třídy
\end{itemize}
\item factory
\begin{itemize}
\item vyrábí instance objektů (typicky podle zadaných parametrů může vyrábět instance různých objektů)
\item nebo může mít factory těch vyráběcích metod více (objekty, které padají z různých metod, spolu typicky nějak souvisí)
\item pokud jsou vyráběcí metody virtuální a můžeme definovat různé factory třídy, které vyrábějí různé typy objektů, tak se tomu říká abstract factory
\end{itemize}
\item visitor
\begin{itemize}
\item uvažujeme potomky A, B, C třídy Base
\item chceme pro každého potomka umět definovat speciální operaci, která se na něm má provést (a pak to chceme umět případně i nějak rozšiřovat)
\item mohli bychom použít \texttt{switch \allowbreak{}(\allowbreak{}obj)\allowbreak{} \allowbreak{}\{ \allowbreak{}case \allowbreak{}Type \allowbreak{}t: \allowbreak{}something(\allowbreak{}t)\allowbreak{};\allowbreak{} \allowbreak{}\}\allowbreak{}}, ale to není objektové ;)
\item místo toho nadefinujeme rozhraní (nebo abstraktní typ) pro visitora, bude obsahovat metody VisitA, VisitB a VisitC (metody se nemusí lišit názvy, stačí, když se liší typem parametrů - jako parametr přijímají ten objekt typu A, B nebo C)
\item třída Base pak bude mít abstraktní metodu \texttt{Accept(\allowbreak{}Visitor \allowbreak{}v)\allowbreak{}}, kterou musí implementovat každý potomek
\item implementace v potomkovi A vypadá tak, že se volá \texttt{v.\allowbreak{}VisitA(\allowbreak{}this)\allowbreak{};\allowbreak{}}
\item použití
\begin{itemize}
\item máme konkrétního visitora \texttt{vis1} a chceme ho použít na objekt \texttt{obj} (který je typu A, B nebo C)
\item zavoláme \texttt{obj.\allowbreak{}Accept(\allowbreak{}vis1)\allowbreak{};\allowbreak{}}
\end{itemize}
\end{itemize}
\end{itemize}
\end{itemize}
\subsection{4. Architektura počítačů a operačních systémů}
\begin{itemize}
\item Základní architektura počítače, reprezentace čísel, dat a programů
\begin{itemize}
\item Harvardská architektura počítače
\begin{itemize}
\item CPU - procesor, vykonává příkazy
\item kódová paměť (code memory) - uložení informací o příkazech
\item komunikační linka - umožňuje, aby procesor četl data z paměti
\begin{itemize}
\item v základním počítači CPU nemusí zapisovat do kódové paměti, do ní se zapisuje nějak z venku
\end{itemize}
\item datová paměť (data memory) - CPU z ní čte a zapisuje do ní data
\begin{itemize}
\item data se ukládají do proměnných
\end{itemize}
\item k procesoru jsou dále připojena vstupně-výstupní zařízení - I/O, periferie
\end{itemize}
\item von Neumannova architektura
\begin{itemize}
\item operační paměť, kde je obojí - v jedné části kód programu, v jiné části data programu
\item zajistíme, aby IP (instruction pointer, jinak také PC, program counter) obsahoval adresy z určitého rozsahu a aby adresy proměnných byly zase z jiného rozsahu
\item je to hardwarově komplikovanější na implementaci, ale přináší to spoustu výhod
\item lze zařídit, aby IP ukazoval na začátek přeloženého programu
\item potenciální zranitelnost, pokud útočník do dat určených ke čtení uloží spustitelné byty a zařídí jejich spuštění
\item operační paměť bude volatile - RAM
\begin{itemize}
\item potřebujeme jiné magické zařízení, které do RAM zkopíruje počáteční program z non-volatile paměti
\end{itemize}
\end{itemize}
\end{itemize}
\item Reprezentace a přístup k datům v paměti, adresa, adresový prostor
\begin{itemize}
\item data i programy reprezentujeme binárně
\item každému bytu v paměti přiřadíme adresu (n-bit unsigned), definujeme si paměťový adresový prostor
\item u 256bytové paměti budeme mít 8bitový adresový prostor (0 až 255 $\dots 2^8-1$)
\item 32bitový adresový prostor stačí pro 4 GB paměti, pro větší paměť už potřebujeme 64 bitů
\end{itemize}
\item Ukládání jednoduchých a složených datových typů
\begin{itemize}
\item jednoduché datové typy
\begin{itemize}
\item čísla se ukládají ve dvojkové soustavě
\item znaménková čísla se ukládají ve dvojkovém doplňku
\begin{itemize}
\item kladná čísla jako unsigned
\item záporná čísla jsou negace absolutní hodnoty plus 1
\item rozsah -128 až 127
\item nefunguje porovnání mezi kladnými a zápornými čísly - procesor musí podporovat znaménkové porovnání
\item MSb určuje znaménko
\item nejběžněji používaná implementace záporných čísel
\end{itemize}
\item čísla s plovoucí desetinnou čárkou (floating-point) se reprezentují takto: $\text{value}=(-1)^{\text{sign}}\cdot\text{significand}\cdot 2^{\text{exponent}-\text{bias}}$
\item číslo se naseká na byty a pak se ukládá v paměti, pořadí uložení bytů je dáno endianitou (little-endian vs. big-endian)
\item text se ukládá pomocí kódování, základní znaky používají ASCII, dále se obvykle používá Unicode
\end{itemize}
\item složené datové typy
\begin{itemize}
\item moderní procesory vyžadují, aby byla data v paměti zarovnaná podle jejich velikosti
\item např. 4bajtový int musí mít adresu zarovnanou na 4 (dělitelnou čtyřmi)
\item celá struktura (struct) je zarovnaná na největší datový typ dostupný na CPU (např. 16B)
\item některé jazyky přehazují položky ve struktuře, aby se eliminovalo volné místo
\item sizeof vrátí velikost včetně těch mezer 
\end{itemize}
\end{itemize}
\item Implementovat běžné programové konstrukce vyšších jazyků (přiřazení, podmínka, cyklus, volání funkce) pomocí instrukcí procesoru
\begin{practicebox}
Překlad konstrukcí (přiřazení, podmínka, cyklus, volání funkce) do instrukcí procesoru je dovednost - natrénuj na konkrétní (i fiktivní) instrukční sadě.
\end{practicebox}
\begin{itemize}
\item nacvičit
\end{itemize}
\item Zapsat běžnou konstrukci vyššího jazyka (přiřazení, podmínka, cyklus, volání funkce), která odpovídá zadané sekvenci (vysvětlených) instrukcí procesoru
\begin{practicebox}
Opačný směr (z dané sekvence instrukcí poznat konstrukci vyššího jazyka) je třeba natrénovat na příkladech.
\end{practicebox}
\begin{itemize}
\item nacvičit
\end{itemize}
\item Podpora pro běh operačního systému - privilegovaný a neprivilegovaný režim procesoru, jádro operačního systému
\begin{itemize}
\item režimy procesoru (bývá jich několik, ale nejdůležitější jsou dva)
\item uživatelský režim (neprivilegovaný)
\begin{itemize}
\item přístupný všem aplikacím
\item omezený přístup ke zdrojům (systémovým registrům, instrukcím)
\end{itemize}
\item kernel mode (privilegovaný)
\begin{itemize}
\item je používán operačním systémem nebo jeho částí
\item má plný přístup ke zdrojům
\end{itemize}
\item přechod mezi režimy je jasně definovaný - je pouze omezené množství způsobů, jak přepnout z uživatelského do kernel režimu (aby bylo možné vyhodnotit, zda je to záměr, nebo chyba)
\begin{itemize}
\item syscall = volání systémového API
\item tenká vrstva nad OS zajišťuje syscally
\end{itemize}
\item např. k vypsání textu do konzole je potřeba, abychom se přepnuli do kernel módu a zpátky (zrychluje se to použitím bufferu)
\item jádro OS zajišťuje inicializaci hardwaru, následně řeší správu prostředků a spouštění/obsluhu programů
\item podle velikosti jádra OS se rozlišují architektury
\begin{itemize}
\item monolitická
\begin{itemize}
\item obslužná vrstva, jejím prostřednictvím se volají interní funkce
\item takhle funguje linuxový kernel
\item zranitelnost - programy v jádře mohou dělat cokoliv
\item původně byly součásti jádra pevně nastavené od instalaci
\begin{itemize}
\item problém - při připojení klávesnice chceme nainstalovat driver zařízení $\to$ dnes lze obsah jádra měnit
\end{itemize}
\end{itemize}
\item vrstvená
\begin{itemize}
\item každá vrstva má svoje rozhraní
\item každá vrstva může používat jenom vrstvu přímo pod ní (její rozhraní)
\end{itemize}
\item mikrokernel
\begin{itemize}
\item většina služeb se z kernel space vystrčí do user space, jsou tam jako jednotlivé moduly, ty komunikují bez sebou a s jádrem
\item je to rozšiřitelné, spolehlivé (jednotlivé služby lze v případě chyby restartovat) a bezpečné, ale poměrně pomalé
\item takhle fungují Windows
\begin{itemize}
\item mají ještě hardwarovou abstrakční vrstvu, takže mikrokernel je přenositelný
\item služby běží v kernel režimu (což je v rozporu s filozofií mikrokernel)
\end{itemize}
\end{itemize}
\end{itemize}
\end{itemize}
\item Popsat roli řadiče zařízení při programem řízené obsluze zařízení (PIO), pro zadané adresy a funkce vstupních a výstupních portů implementovat programem řízenou obsluhu zadaného zařízení (myš, disk)
\begin{itemize}
\item implementaci nacvičit
\item řadič zařízení (controller) - HW součástka, komunikuje se zařízením nějakým protokolem
\item ovladač zařízení (driver) - SW komponenta, součást operačního systému, realizuje komunikaci s řadičem, poskytuje operačnímu systému jednotné komunikační rozhraní
\item operační systém pak zařízení zprostředkovává běžícímu programu/procesu (a uživateli)
\end{itemize}
\item Popsat roli přerušení při programem řízené obsluze zařízení (PIO), na úrovni vykonávání instrukcí popsat reakci procesoru (hardware) a operačního systému (software) na žádost o přerušení
\begin{itemize}
\item komunikace se zařízeními
\begin{itemize}
\item polling - CPU se aktivně ptá na změnu stavu zařízení (pomalé, dnes už se nepoužívá)
\item přerušení - řadič vyšle signál, že je operace hotová, a přeruší chod procesoru (ale pořád musím kopírovat data, což je pomalé)
\item DMA (Direct Memory Access) - přenos dat ze zařízení do paměti probíhá hardwarově bez účasti procesoru, až pak se provede přerušení; funkce scatter/gather umožňuje využívat nesouvislé části paměti
\end{itemize}
\item typy přerušení
\begin{itemize}
\item externí - hardwarové pomocí IRQ pinu, lze ho zamaskovat (deaktivovat - pomocí registru)
\item (hardwarová) výjimka - nastal problém s instrukcemi $\to$ procesor vyvolá přerušení a nechá na OS, aby situaci vyřešil
\begin{itemize}
\item výjimka typu trap se volá po instrukci, např. dělení nulou
\item u výjimek typu fault se procesor vrátí na stav před instrukcí, nechám OS, aby problém opravil, a potom instrukci pustím znova
\end{itemize}
\item softwarové - speciální instrukce
\end{itemize}
\item jak funguje přerušení
\begin{itemize}
\item CPU zjistí zdroj přerušení, získá adresu handleru (kdo přerušení zpracuje - je fixní nebo se použije tabulka)
\item proud instrukcí je přerušen, CPU začne vykonávat instrukce handleru (obvykle se přepne do privilegovaného režimu, protože handler je součástí kernelu)
\item handler uloží stav CPU
\item handler udělá něco užitečného
\item handler obnoví stav CPU
\item CPU pokračuje v původním proudu instrukcí
\end{itemize}
\end{itemize}
\item Neprivilegované (uživatelské) procesy
\begin{itemize}
\item k prostředkům přistupují skrze operační systém (ten pak využívá ovladače apod.)
\item mají k dispozici virtuální paměť
\end{itemize}
\item Procesy, vlákna, kontext procesu a vlákna
\begin{itemize}
\item program
\begin{itemize}
\item pasivní soubor na disku
\item loader - vezme program a načte ho do paměti
\item někde v hlavičce programu je informace o tom, kde program (výkon instrukcí) začíná
\end{itemize}
\item proces
\begin{itemize}
\item instance programu vytvořená operačním systémem
\item je to datová struktura uvnitř operačního systému, v níž jsou uložené různá data, která daný program potřebuje
\item vlastní: kód, prostor v paměti, další prostředky
\end{itemize}
\item vlákno (thread)
\begin{itemize}
\item místo uvnitř procesu, kde se vykonávají instrukce
\item kontext procesoru - datová struktura, ve které jsou uložené registry procesoru, pokud dané vlákno zrovna neběží
\item vlastní: pozici v kódu (program counter / instruction pointer), svůj zásobník, stav procesoru
\end{itemize}
\item fiber
\begin{itemize}
\item lightweight vlákno
\item nemá celý kontext
\item scheduling se dělá kooperativně (jednotlivé fibery se vzájemně vyměňují v běhu)
\end{itemize}
\end{itemize}
\item Přepínání kontextu, kooperativní a preemptivní multitasking
\begin{itemize}
\item scheduler - část operačního systému, používá schedulingové algoritmy k~přidělení zdrojů (jader) jednotkám
\item když se přeruší vykonávání vlákna, provede se context switch - kontext procesoru (registry včetně PC) se uloží
\item multitasking
\begin{itemize}
\item cooperative - všechny procesy kooperují, předávají si řízení
\item preemptive - každé vlákno dostane od scheduleru svůj time slice; jakmile čas vyprší, dojde k přerušení
\end{itemize}
\item jinými slovy
\begin{itemize}
\item při kooperativním multitaskingu musí proces dobrovolně předat řízení (yield)
\item při preemptivním multitaskingu plánovač proces prostě přeruší, jakmile mu dojde čas
\end{itemize}
\end{itemize}
\item Plánování běhu procesů a vláken, stavy vlákna
\begin{itemize}
\item stavy vlákna
\begin{itemize}
\item created
\item ready
\item running
\item blocked
\item terminated (zombie)
\end{itemize}
\item scheduler - část operačního systému, používá schedulingové algoritmy k~přidělení zdrojů (jader) jednotkám
\item real-time scheduling: real-time proces má čas, kdy se má spustit, a čas, do kdy má skončit (hard deadline - nemá smysl pokračovat, soft deadline - má smysl pokračovat ve výpočtu)
\item priorita - vázaná na proces
\begin{itemize}
\item je dvousložková - při spuštění se přiděluje statická priorita (podle uživatele), dynamická priorita se v časových intervalech pravidelně zvyšuje všem ready vláknům (po spuštění se dynamická priorita sníží na nula); celková priorita je daná součtem
\end{itemize}
\item kooperativní algoritmy
\begin{itemize}
\item first come, first serve (FCFS)
\begin{itemize}
\item FIFO řada, procesy se vykonávají v pořadí, ve kterém přišly
\end{itemize}
\item shortest job first
\begin{itemize}
\item musí být známý očekávaný čas běhu
\end{itemize}
\item longest job first
\end{itemize}
\item preemptivní algoritmy
\begin{itemize}
\item round robin
\begin{itemize}
\item jako FCFS
\item když vlákno trvá moc dlouho, tak se zařadí na konec fronty
\end{itemize}
\item multilevel feedback-queue
\begin{itemize}
\item každá fronta má jinak dlouhý time slice
\item první (horní) ho má kratší, ty pod ní ho mají delší
\item když vlákno trvá moc dlouho, tak se zařadí na konec nižší fronty
\item když se vlákno brzo zablokuje, tak se zařadí na konec vyšší fronty
\end{itemize}
\item completely fair scheduler (CFS)
\begin{itemize}
\item používá červeno-černý strom
\end{itemize}
\end{itemize}
\end{itemize}
\item Vysvětlit rozdíl mezi virtuální a fyzickou adresou a identifikovat, zda se v zadaném kontextu či fragmentu kódu používá virtuální nebo fyzická adresa
\begin{itemize}
\item instrukce procesoru pracují s virtuálními adresami
\item operační paměť má fyzické adresy
\item hardwarově je implementován převodní mechanismus
\begin{itemize}
\item musí vyhodnotit, jestli existuje převod
\item zajistí přepočet, pokud existuje
\item tyto kroky musí být extrémně rychlé
\end{itemize}
\item existují dva mechanismy - segmentace (už se nepoužívá) a stránkování
\item MMU = Memory Management Unit, provádí převod
\item převod paměti je transparentní
\item výhody konceptu
\begin{itemize}
\item větší adresový prostor - abych mohl spustit proces s většími požadavky na paměť
\item bezpečnost - aby si procesy navzájem nemohly koukat do paměti, každý má vlastní mapování; to je dnes ten hlavní důvod
\end{itemize}
\end{itemize}
\item Na zadaném příkladu identifikovat a vysvětlit význam komponent virtuální a fyzické adresy (číslo stránky, číslo rámce, offset)
\begin{itemize}
\item virtuální adresový prostor (VAS) je rozdělen na stejné části (stránky, jejich velikost odpovídá mocninám dvojky)
\item fyzický adresový prostor (PAS) je rozdělen na stejné části (rámce, framy, mají stejnou velikost jako stránky)
\item VAS je jednorozměrný, instrukce pracují s jedním číslem
\item page table (stránkovací tabulka)
\begin{itemize}
\item zajišťuje překlad adres
\item je to pole stránek (indexované číslem stránky)
\item každý záznam obsahuje číslo framu a atributy
\item je tam jednička/nula - aby se dalo určit, jestli stránka fyzicky existuje
\item v případě chyby - page fault
\end{itemize}
\item virtuální adresa se navenek jeví jako jedno číslo, ale protože velikost stránky je mocnina dvojky, tak část binárního čísla odpovídá číslu stránky a část offsetu
\begin{itemize}
\item offset je u virtuální i fyzické adresy stejný, tedy se překládá jenom prefix odpovídající stránce (přeloží se na číslo rámce)
\end{itemize}
\item když mi dojde fyzická paměť, tak nějakou stránku uložím na disk - je to obvykle méně dat, než by to bylo při výpadku segmentu (segmenty měly různé velikosti, stránky jsou všechny stejně velké)
\end{itemize}
\item Pro konkrétní adresy a obsah jednoúrovňové stránkovací tabulky řešit úlohy překladu adres
\begin{itemize}
\item ukázkový příklad:
\begin{itemize}
\item 32bitová virtuální adresa
\item posledních 12 bitů určuje offset
\item prvních 20 bitů určuje stránku, použiju je k indexaci do stránkovací tabulky
\item dozvím se číslo rámce dat
\item k číslu rámce připojím offset
\item dostanu 32bitovou fyzickou adresu
\end{itemize}
\end{itemize}
\item Vysvětlit roli virtuálních adresových prostorů v ochraně paměti procesů a vláken
\begin{itemize}
\item proces má vlastní virtuální adresový prostor, takže se nemůže stát, že by sahal na paměť ostatních procesů
\item ale vlákna stejného procesu virtuální adresový prostor sdílejí
\item někdy naopak chceme, aby různé procesy mohly sdílet nějaké místo v paměti - dá se zařídit, aby konkrétní virtuální adresy různých procesů ukazovaly na stejné místo do fyzické paměti
\end{itemize}
\item Sdílení úložného prostoru - soubory, analogie s adresovým prostorem; abstrakce a rozhraní pro práci se soubory
\begin{itemize}
\item soubor
\begin{itemize}
\item jednotka organizace dat, kolekce souvisejících informací
\item jádro OS nerozumí formátům souborů
\item soubory se typicky neukládají v paměti, ale na perzistentním úložišti
\item soubor je chápán jako unikátní číslo - kvůli lidem mají soubory jméno a cestu (aby v souborech nebyl chaos)
\begin{itemize}
\item analogie s adresovým prostorem - cestu souboru je nutné přeložit (jako se virtuální adresa překládá na fyzickou)
\end{itemize}
\item některé části názvu souboru (počáteční tečka, přípona) mají speciální význam (ale přípona souboru nemusí souviset s reálným formátem souboru)
\item operace se soubory - vyčištění cache, změna atributů, vytvoření, smazání
\item file handle - číslo specifické pro proces, odkazuje na konkrétní soubor (stdin, stdout, stderr mají handles 0, 1, 2)
\end{itemize}
\item buffering
\begin{itemize}
\item cachování sektorů disku (když disk čtu po bajtech, první se načítá přímo z disku, všechny ostatní v daném sektoru už se pak načítají z cache)
\item existuje několik úrovní cache (systémová, runtime)
\item sekvenční vs. náhodný přístup
\end{itemize}
\item alternativy - memory mapping, asynchronní přístup souborům
\item adresář
\begin{itemize}
\item kolekce souborů
\item význam - rychlejší hledání souboru, lepší navigace pro uživatele, logické seskupení souborů
\item obvykle reprezentován jako soubor
\item někdy v něm můžou být uloženy některé atributy jeho souborů
\item obvykle existuje kořenový adresář
\item operace - vytvořit, smazat, přejmenovat, najít název, vypsat soubory
\end{itemize}
\item ukládání souborů
\begin{itemize}
\item tradiční úložiště
\begin{itemize}
\item na sekundárním nebo externím úložišti
\item file system in RAM (pro dočasné soubory)
\end{itemize}
\item síťové úložiště - soubory se tváří, jako by byly na disku, ale přitom jsou uloženy někde jinde na síti
\item virtuální soubory - např. /dev/null
\end{itemize}
\item file links
\begin{itemize}
\item links (hard links)
\begin{itemize}
\item více adresářových záznamů odkazuje na stejný fyzický soubor
\item většina operací je transparentních
\item šetří místo, ale může dělat problémy (smazání hardlinku nesmí smazat fyzický soubor - takže OS musí počítat odkazy na daný soubor)
\end{itemize}
\item symlinks (soft links)
\begin{itemize}
\item symlinky jsou speciální soubory, v nichž je uložená cesta jiného souboru
\end{itemize}
\end{itemize}
\item file system
\begin{itemize}
\item řeší překlad jmen, správu datových bloků, správu dat souborů
\item pro file system je jednotkou jeden blok, blok je určitý počet sektorů
\item řeší abstrakci práce se soubory a adresáři
\item může být lokální (FAT, NTFS) nebo síťový (NFS, CIFS/SMB)
\end{itemize}
\end{itemize}
\item Časově závislé chyby (race conditions)
\begin{itemize}
\item více vláken přistupuje ke sdílenému prostředku
\item to vede k tomu, že výsledek výpočtu závisí na plánování operačního systému nebo na chování procesoru
\item takový výsledek je k ničemu
\item řešila by to atomizace read-modify-write operace
\item jak to vyřešit?
\begin{itemize}
\item definuju kritickou sekci
\item pomocí synchronizačního primitiva zajistím, že v kritické sekci je jenom jedno vlákno
\end{itemize}
\end{itemize}
\item Kritická sekce, vzájemné vyloučení
\begin{itemize}
\item kritická sekce - (nejmenší) kus zdrojového kódu, kde dochází k přístupu ke sdílenému prostředku, který nesmí používat víc procesů najedou
\item vzájemné vyloučení (mutex) - jednoduchý algoritmus (synchronizační primitivum) zajišťující, že do kritické sekce nevstoupí víc vláken/procesů najednou
\begin{itemize}
\item je to vlastně zámek (ale zámek se obvykle týká vláken, kdežto mutex se vztahuje na procesy)
\item pokud je mutex zamčený nějakým procesem, jiný proces nemůže sdílený prostředek používat
\item umožňuje vícenásobné zamčení tím stejným procesem (pak je nezbytný stejný počet odemčení)
\end{itemize}
\end{itemize}
\item Základní synchronizační primitiva, jejich rozhraní a použití - aktivní a pasivní čekání, zámky
\begin{itemize}
\item aktivní a pasivní/blokující synchronizační primitiva
\begin{itemize}
\item aktivní vykonávají instrukce a pořád koukají do kritické sekce, jestli tam můžou
\item pasivní/blokující jsou zablokovány, dokud není přístup povolen
\end{itemize}
\item hardwarová podpora - atomické instrukce test-and-set (TSL), compare-and-swap (CAS)
\item primitiva
\begin{itemize}
\item spin-lock - aktivní čekání pomocí TSL/CAS, ideální pro krátké čekání
\item semafor - counter a fronta čekajících vláken, atomické operace UP a DOWN
\item mutex - semafor s počítadlem rovným jedné (umožňuje vstup jednoho vlákna do kritické sekce v danou chvíli), atomické operace LOCK a UNLOCK
\item barrier - několik vláken se v jednu chvíli setká u stejné bariéry
\end{itemize}
\end{itemize}
\end{itemize}
